\chapter{svn Reference---Subversion Command-Line
Client}\label{svn.ref.svn}

\texttt{svn} is the official command-line client of Subversion. Its
functionality is offered via a collection of task-specific subcommands,
most of which accept a number of options for fine-grained control of the
program\textquotesingle s behavior.

When using the \texttt{svn} program, subcommands and other non-option
arguments must appear in a specified order on the command line. Options,
on the other hand, may appear anywhere on the command line (after the
program name, of course), and in general, their order is irrelevant. For
example, all of the following are valid ways to use
\texttt{svn\ status}, and are interpreted in exactly the same way:

\begin{verbatim}
$ svn -vq status myfile
$ svn status -v -q myfile
$ svn -q status -v myfile
$ svn status -vq myfile
$ svn status myfile -qv
\end{verbatim}

The following sections describe each of the various subcommands and
options provided by the \texttt{svn} command-line client program,
including some examples of each subcommand\textquotesingle s typical
uses.

While Subversion has different options for its subcommands, all options
exist in a single namespace---that is, each option is guaranteed to mean
the roughly same thing regardless of the subcommand you use it with. For
example, \texttt{-\/-verbose} (\texttt{-v}) always means ``verbose
output,'' regardless of the subcommand you use it with.

The \texttt{svn} command-line client usually exits quickly with an error
if you pass it an option which does not apply to the specified
subcommand. But as of Subversion 1.5, several of the options which apply
to all---or nearly all---of the subcommands have been deemed acceptable
by all subcommands, even if they have no effect on some of them. (This
change was made primarily to improve the client\textquotesingle s
ability to called from custom wrapping scripts.) These options appear
grouped together in the command-line client\textquotesingle s usage
messages as global options, as can be seen in the following bit of
output:

\begin{verbatim}
$ svn help upgrade
upgrade: Upgrade the metadata storage format for a working copy.
usage: upgrade [WCPATH...]

  Local modifications are preserved.

Valid options:
  -q [--quiet]             : print nothing, or only summary information

Global options:
  --username ARG           : specify a username ARG
  --password ARG           : specify a password ARG
  --no-auth-cache          : do not cache authentication tokens
  --non-interactive        : do no interactive prompting (default is to prompt
                             only if standard input is a terminal device)
  --force-interactive      : do interactive prompting even if standard input
                             is not a terminal device
  --trust-server-cert      : accept SSL server certificates from unknown
                             certificate authorities without prompting (but only
                             with '--non-interactive')
  --config-dir ARG         : read user configuration files from directory ARG
  --config-option ARG      : set user configuration option in the format:
                                 FILE:SECTION:OPTION=[VALUE]
                             For example:
                                 servers:global:http-library=serf
$
\end{verbatim}

\texttt{svn} subcommands recognize the following options:

\protect\phantomsection\label{svn.ref.svn.sw}{}

\begin{description}
\item[\protect\phantomsection\label{svn.ref.svn.sw.accept}{}\texttt{-\/-accept}
\textless ACTION\textgreater{}]
Specifies an action for automatic conflict resolution, disabling the
interactive prompts which ask the user how to handle each conflict as it
is noticed. Though which of the specific actions are applicable differs
depending on which subcommand is in use, Subversion supports the
following long (and short) values for \textless ACTION\textgreater:

\begin{description}
\item[\texttt{postpone} (\texttt{p})]
Take no resolution action at all and instead allow the conflicts to be
recorded for future resolution.
\item[\texttt{edit} (\texttt{e})]
Open each conflicted file in a text editor for manual resolution of
line-based conflicts.
\item[\texttt{launch} (\texttt{l})]
Launch an interactive merge conflict resolution tool for each conflicted
file.
\item[\texttt{base}]
Choose the file that was the (unmodified) \texttt{BASE} revision before
you tried to integrate changes from the server into your working copy.
\item[\texttt{working}]
Assuming that you\textquotesingle ve manually handled the conflict
resolution, choose the version of the file as it currently stands in
your working copy.
\item[\texttt{mine-full} (\texttt{mf})]
Resolve conflicted files by preserving all local modifications and
discarding all changes fetched from the server during the operation
which caused the conflict.
\item[\texttt{theirs-full} (\texttt{tf})]
Resolve conflicted files by discarding all local modifications and
integrating all changes fetched from the server during the operation
which caused the conflict.
\item[\texttt{mine-conflict} (\texttt{mc})]
Resolve conflicted files by preferring local modifications over the
changes fetched from the server in conflicting regions of each
file\textquotesingle s content.
\item[\texttt{theirs-conflict} (\texttt{tc})]
Resolve conflicted files by preferring the changes fetched from the
server over local modifications in conflicting regions of each
file\textquotesingle s content.
\end{description}

Consult the output of \texttt{svn\ help\ SUBCOMMAND} to see exactly
which actions are supported by the specific subcommand of interest.
\item[\protect\phantomsection\label{svn.ref.svn.sw.allow_mixed_revisions}{}\texttt{-\/-allow-mixed-revisions}]
Disables the verification---performed by default by \texttt{svn\ merge}
as of Subversion 1.7---that the target of a merge operation and all of
its children are at a uniform revision. While merging into a
single-revision working copy target is the recommended best practice,
this option may be used to permit merges into mixed-revision working
copies as necessary.
\item[\protect\phantomsection\label{svn.ref.svn.sw.auto_props}{}\texttt{-\/-auto-props}]
Enables automatic property assignment (per runtime configuration rules),
overriding the \texttt{enable-auto-props} runtime configuration
directive.
\item[\protect\phantomsection\label{svn.ref.svn.sw.change}{}\texttt{-\/-change}
(\texttt{-c}) \textless ARG\textgreater{}]
Perform the requested operation using a specific ``change''. Generally
speaking, this option is syntactic sugar for \texttt{-r\ ARG-1:ARG}.
Some subcommands permit a comma-separated list of revision number
arguments (e.g., \texttt{-c\ ARG1,ARG2,ARG3}). Alternatively, you can
provide two arguments separated by a dash (as in \texttt{-c\ ARG1-ARG2})
to identify the range of revisions between \textless ARG1\textgreater{}
and \textless ARG2\textgreater, inclusive. Finally, if the revision
argument is negated, the implied revision range is reversed:
\texttt{-c\ -45} is equivalent to \texttt{-r\ 45:44}.
\item[\protect\phantomsection\label{svn.ref.svn.sw.changelist}{}\texttt{-\/-changelist}
(\texttt{-\/-cl}) \textless ARG\textgreater{}]
Instructs Subversion to operate only on members of the changelist named
\textless ARG\textgreater. You can use this option multiple times to
specify sets of changelists.
\item[\protect\phantomsection\label{svn.ref.svn.sw.config_dir}{}\texttt{-\/-config-dir}
\textless DIR\textgreater{}]
Instructs Subversion to read configuration information from the
specified directory instead of the default location
(\texttt{.subversion} in the user\textquotesingle s home directory).

This option is accepted by all \texttt{svn} subcommands.
\item[\protect\phantomsection\label{svn.ref.svn.sw.config_option}{}\texttt{-\/-config-option}
\textless CONFSPEC\textgreater{}]
Sets, for the duration of the command, the value of a runtime
configuration option. \textless CONFSPEC\textgreater{} is a string which
specifies the configuration option namespace, name and value that
you\textquotesingle d like to assign, formatted as
\textless FILE\textgreater:\textless SECTION\textgreater:\textless OPTION\textgreater={[}\textless VALUE\textgreater{]}.
In this syntax, \textless FILE\textgreater{} and
\textless SECTION\textgreater{} are the runtime configuration file
(either \texttt{config} or \texttt{servers}) and the section thereof,
respectively, which contain the option whose value you wish to change.
\textless OPTION\textgreater{} is, of course, the option itself, and
\textless VALUE\textgreater{} the value (if any) you wish to assign to
the option. For example, to temporarily disable the use of the
compression in the HTTP protocol, use
\texttt{-\/-config-option=servers:global:http-compression=no}. You can
use this option multiple times to change multiple option values
simultaneously.

This option is accepted by all \texttt{svn} subcommands.
\item[\protect\phantomsection\label{svn.ref.svn.sw.depth}{}\texttt{-\/-depth}
\textless ARG\textgreater{}]
Instructs Subversion to limit the scope of an operation to a particular
tree depth. \textless ARG\textgreater{} is one of \texttt{empty} (only
the target itself), \texttt{files} (the target and any immediate file
children thereof), \texttt{immediates} (the target and any immediate
children thereof), or \texttt{infinity} (the target and all of its
descendants---full recursion).
\item[\protect\phantomsection\label{svn.ref.svn.sw.diff}{}\texttt{-\/-diff}]
Enables a special output mode for \texttt{svn\ log} which includes a
difference report (a la \texttt{svn\ diff}) as part of each
revision\textquotesingle s information.
\item[\protect\phantomsection\label{svn.ref.svn.sw.diff_cmd}{}\texttt{-\/-diff-cmd}
\textless CMD\textgreater{}]
Specifies an external program to use to show differences between files.
When \texttt{svn\ diff} is invoked without this option, it uses
Subversion\textquotesingle s internal differencing engine, which
provides unified diffs by default. If you want to use an external
differencing program, use \texttt{-\/-diff-cmd}. You can then pass
options to the specified program using the \texttt{-\/-extensions}
(\texttt{-x}) option.
\item[\protect\phantomsection\label{svn.ref.svn.sw.diff3_cmd}{}\texttt{-\/-diff3-cmd}
\textless CMD\textgreater{}]
Specifies an external 3-way differencing program (used to merge
line-based changes into files).
\item[\protect\phantomsection\label{svn.ref.svn.sw.dry_run}{}\texttt{-\/-dry-run}]
Goes through all the motions of running a command, but makes no actual
changes---either on disk or in the repository.
\item[\protect\phantomsection\label{svn.ref.svn.sw.editor_cmd}{}\texttt{-\/-editor-cmd}
\textless CMD\textgreater{}]
Specifies an external program to use to edit a log message or a property
value. See the \texttt{editor-cmd} section in
\hyperref[svn.advanced.confarea.opts.config]{General configuration} for
ways to specify a default editor.
\item[\protect\phantomsection\label{svn.ref.svn.sw.encoding}{}\texttt{-\/-encoding}
\textless ENC\textgreater{}]
Tells Subversion that your commit message is composed using the
character encoding provided. The default character encoding is derived
from your operating system\textquotesingle s native locale; use this
option if your commit message is composed using any other encoding.
\item[\protect\phantomsection\label{svn.ref.svn.sw.extensions}{}\texttt{-\/-extensions}
(\texttt{-x}) \textless ARG\textgreater{}]
Specifies customizations which Subversion should make when performing
difference calculations. Valid extensions include:

\begin{description}
\item[\texttt{-\/-ignore-space-change} (\texttt{-b})]
Ignore changes in the amount of white space.
\item[\texttt{-\/-ignore-all-space} (\texttt{-w})]
Ignore all white space.
\item[\texttt{-\/-ignore-eol-style}]
Ignore changes in EOL (end-of-line) style.
\item[\texttt{-\/-show-c-function} (\texttt{-p})]
Show C function names in the diff output.
\item[\texttt{-\/-unified} (\texttt{-u})]
Show three lines of unified diff context.
\end{description}

The default value of \textless ARG\textgreater{} is \texttt{-u}. If you
wish to pass multiple arguments, you must enclose all of them in quotes.

Note that when Subversion is configured to invoke an external diff
command, the value of the \texttt{-\/-extensions} (\texttt{-x}) option
isn\textquotesingle t restricted to the previously mentioned options,
but may be \emph{any} additional arguments which Subversion should pass
to that command.
\item[\protect\phantomsection\label{svn.ref.svn.sw.file}{}\texttt{-\/-file}
(\texttt{-F}) \textless FILENAME\textgreater{}]
Uses the contents of the named file for the specified subcommand.
Different subcommands do different things with this content. For
example, \texttt{svn\ commit} uses the content as a commit log message,
whereas \texttt{svn\ propset} uses it as a property value.
\item[\protect\phantomsection\label{svn.ref.svn.sw.force}{}\texttt{-\/-force}]
Forces a particular command or operation to run. Subversion will prevent
you from performing some operations in normal usage, but you can pass
this option to tell Subversion ``I know what I\textquotesingle m doing
as well as the possible repercussions of doing it, so let me at
\textquotesingle em.'' This option is the programmatic equivalent of
doing your own electrical work with the power on---if you
don\textquotesingle t know what you\textquotesingle re doing,
you\textquotesingle re likely to get a nasty shock.
\item[\protect\phantomsection\label{svn.ref.svn.sw.force_log}{}\texttt{-\/-force-log}]
Forces a suspicious parameter passed to the \texttt{-\/-message}
(\texttt{-m}) or \texttt{-\/-file} (\texttt{-F}) option to be accepted
as valid. By default, Subversion will produce an error if parameters to
these options look like they might instead be targets of the subcommand.
For example, if you pass a versioned file\textquotesingle s path to the
\texttt{-\/-file} (\texttt{-F}) option, Subversion will assume
you\textquotesingle ve made a mistake, that the path was instead
intended as the target of the operation, and that you simply failed to
provide some other---unversioned---file as the source of your log
message. To assert your intent and override these types of errors, pass
the \texttt{-\/-force-log} option to subcommands that accept log
messages.
\item[\protect\phantomsection\label{svn.ref.svn.sw.force_interactive}{}\texttt{-\/-force-interactive}]
Forces the \texttt{svn} command-line client to run in interactive mode
when standard input is not a terminal device.

This option is accepted by all \texttt{svn} subcommands.
\item[\protect\phantomsection\label{svn.ref.svn.sw.git}{}\texttt{-\/-git}]
Enables a special output mode for \texttt{svn\ diff} designed for
cross-compatibility with the popular Git distributed version control
system.
\item[\protect\phantomsection\label{svn.ref.svn.sw.help}{}\texttt{-\/-help}
(\texttt{-h}, \texttt{-?})]
If used with one or more subcommands, shows the built-in help text for
each. If used alone, it displays the general client help text.
\item[\protect\phantomsection\label{svn.ref.svn.sw.ignore_ancestry}{}\texttt{-\/-ignore-ancestry}]
Tells Subversion to ignore ancestry when calculating differences (rely
on path contents alone). Also disables
\hyperref[svn.branchmerge.basicmerging.mergetracking]{Merge Tracking}
when used with the \texttt{svn\ merge} subcommand.
\item[\protect\phantomsection\label{svn.ref.svn.sw.ignore_externals}{}\texttt{-\/-ignore-externals}]
Tells Subversion to ignore externals definitions and the external
working copies managed by them.
\item[\protect\phantomsection\label{svn.ref.svn.sw.ignore_keywords}{}\texttt{-\/-ignore-keywords}]
Disables keyword expansion.
\item[\protect\phantomsection\label{svn.ref.svn.sw.ignore_properties}{}\texttt{-\/-ignore-properties}]
Instructs \texttt{svn\ diff} to suppress output of property changes.
\item[\protect\phantomsection\label{svn.ref.svn.sw.ignore_whitespace}{}\texttt{-\/-ignore-whitespace}]
Instructs \texttt{svn\ patch} to ignore whitespace when attempting to
identify patch context.
\item[\protect\phantomsection\label{svn.ref.svn.sw.incremental}{}\texttt{-\/-incremental}]
Prints output in a format suitable for concatenation to prior similar
output.
\item[\protect\phantomsection\label{svn.ref.svn.sw.internal_diff}{}\texttt{-\/-internal-diff}]
Instructs Subversion to use its built-in differencing engine despite any
external differencing mechanism that may be specified for use in the
user\textquotesingle s runtime configuration.
\item[\protect\phantomsection\label{svn.ref.svn.sw.keep_changelists}{}\texttt{-\/-keep-changelists}]
Tells Subversion not to remove the changelist assigments from working
copy items after committing.
\item[\protect\phantomsection\label{svn.ref.svn.sw.keep_local}{}\texttt{-\/-keep-local}]
Keeps the local copy of a file or directory (used with the
\texttt{svn\ delete} command).
\item[\protect\phantomsection\label{svn.ref.svn.sw.limit}{}\texttt{-\/-limit}
(\texttt{-l}) \textless NUM\textgreater{}]
Shows only the first \textless NUM\textgreater{} log messages.
\item[\protect\phantomsection\label{svn.ref.svn.sw.message}{}\texttt{-\/-message}
(\texttt{-m}) \textless MESSAGE\textgreater{}]
Indicates that you will specify either a log message or a lock comment
on the command line, following this option. For example:

\begin{verbatim}
$ svn commit -m "They don't make Sunday."
\end{verbatim}
\item[\protect\phantomsection\label{svn.ref.svn.sw.native_eol}{}\texttt{-\/-native-eol}
\textless ARG\textgreater{}]
Causes \texttt{svn\ export} to use a specific end-of-line sequence as if
it was the native sequence for the client platform.
\textless ARG\textgreater{} may be one of \texttt{CR}, \texttt{LF}, or
\texttt{CRLF}.
\item[\protect\phantomsection\label{svn.ref.svn.sw.new}{}\texttt{-\/-new}
\textless ARG\textgreater{}]
Uses \textless ARG\textgreater{} as the newer target (for use with
\texttt{svn\ diff}).
\item[\protect\phantomsection\label{svn.ref.svn.sw.no_auth_cache}{}\texttt{-\/-no-auth-cache}]
Prevents caching of authentication information (e.g., username and
password) in the Subversion runtime configuration directories.

This option is accepted by all \texttt{svn} subcommands.
\item[\protect\phantomsection\label{svn.ref.svn.sw.no_auto_props}{}\texttt{-\/-no-auto-props}]
Disables automatic property setting, overriding the
\texttt{enable-auto-props} runtime configuration directive.
\item[\protect\phantomsection\label{svn.ref.svn.sw.no_diff_added}{}\texttt{-\/-no-diff-added}]
Prevents Subversion from printing differences for added files. The
default behavior when you add a file is to print the same differences
that you would see if you had added the entire contents of an existing
(empty) file.
\item[\protect\phantomsection\label{svn.ref.svn.sw.no_diff_deleted}{}\texttt{-\/-no-diff-deleted}]
Prevents Subversion from printing differences for deleted files. The
default behavior when you remove a file is for \texttt{svn\ diff} to
print the same differences that you would see if you had kept the file
but removed all of its content.
\item[\protect\phantomsection\label{svn.ref.svn.sw.no_ignore}{}\texttt{-\/-no-ignore}]
Shows files in the status listing or adds/imports files that would
normally be omitted since they match a pattern in the
\texttt{global-ignores} configuration option or the \texttt{svn:ignore}
or \texttt{svn:global-ignores}properties. See
\hyperref[svn.advanced.confarea.opts.config]{General configuration} and
\hyperref[svn.advanced.props.special.ignore]{Ignoring Unversioned Items}
for more information.
\item[\protect\phantomsection\label{svn.ref.svn.sw.no_unlock}{}\texttt{-\/-no-unlock}]
Tells Subversion not to automatically unlock files. (The default commit
behavior is to unlock all files listed as part of the commit.) See
\hyperref[svn.advanced.locking]{Locking} for more information.
\item[\protect\phantomsection\label{svn.ref.svn.sw.non_interactive}{}\texttt{-\/-non-interactive}]
Disables all interactive prompting. Some examples of interactive
prompting include requests for authentication credentials and conflict
resolution decisions. This is useful if you\textquotesingle re running
Subversion inside an automated script and it\textquotesingle s more
appropriate to have Subversion fail than to prompt for more information.

Beginning with Subversion 1.8, the \texttt{svn} command-line client, by
default, is non-interactive when standard input is not a terminal
device. Pass the \texttt{-\/-force-interactive} option to make the
client run in interactive mode.

This option is accepted by all \texttt{svn} subcommands.
\item[\protect\phantomsection\label{svn.ref.svn.sw.non_recursive}{}\texttt{-\/-non-recursive}
(\texttt{-N})]
\emph{Deprecated}. Stops a subcommand from recursing into
subdirectories. Most subcommands recurse by default, but some do not.
Users should avoid this option and use the more precise
\texttt{-\/-depth} option instead. For most subcommands, specifying
\texttt{-\/-non-recursive} produces behavior which is the same as if
you\textquotesingle d specified \texttt{-\/-depth=files}, but there are
exceptions: non-recursive \texttt{svn\ status} operates at the
\texttt{immediates} depth, and the non-recursive forms of
\texttt{svn\ revert}, \texttt{svn\ add}, and \texttt{svn\ commit}
operate at an \texttt{empty} depth.
\item[\protect\phantomsection\label{svn.ref.svn.sw.notice_ancestry}{}\texttt{-\/-notice-ancestry}]
Pays attention to ancestry when calculating differences.
\item[\protect\phantomsection\label{svn.ref.svn.sw.old}{}\texttt{-\/-old}
\textless ARG\textgreater{}]
Uses \textless ARG\textgreater{} as the older target (for use with
\texttt{svn\ diff}).
\item[\protect\phantomsection\label{svn.ref.svn.sw.parents}{}\texttt{-\/-parents}]
Creates and adds nonexistent or nonversioned parent directories to the
working copy or repository as part of an operation. This is useful for
automatically creating multiple subdirectories where none currently
exist. If performed on a URL, all the directories will be created in a
single commit.
\item[\protect\phantomsection\label{svn.ref.svn.sw.password}{}\texttt{-\/-password}
\textless PASSWD\textgreater{}]
Specifies the password to use when authenticating against a Subversion
server. If not provided, or if incorrect, Subversion will prompt you for
this information as needed.

This option is accepted by all \texttt{svn} subcommands.
\item[\protect\phantomsection\label{svn.ref.svn.sw.patch_compatible}{}\texttt{-\/-patch-compatible}]
Instructs \texttt{svn\ diff} to produce output compatible with generic
third-party patch tools. The result of using this option is the same as
running \texttt{svn\ diff} with
\texttt{-\/-show-copies-as-adds\ -\/-ignore-properties} options.
\item[\protect\phantomsection\label{svn.ref.svn.sw.properties_only}{}\texttt{-\/-properties-only}]
Instructs \texttt{svn\ diff} to show only property changes.
\item[\protect\phantomsection\label{svn.ref.svn.sw.quiet}{}\texttt{-\/-quiet}
(\texttt{-q})]
Requests that the client print only essential information while
performing an operation.
\item[\protect\phantomsection\label{svn.ref.svn.sw.record_only}{}\texttt{-\/-record-only}]
Enables a special mode of \texttt{svn\ merge} in which the specified
merge operation is recorded in the local merge tracking information, but
is not actually performed.
\item[\protect\phantomsection\label{svn.ref.svn.sw.recursive}{}\texttt{-\/-recursive}
(\texttt{-R})]
Makes a subcommand recurse into subdirectories. (Most subcommands
recurse by default.)
\item[\protect\phantomsection\label{svn.ref.svn.sw.reintegrate}{}\texttt{-\/-reintegrate}]
\emph{Deprecated}. Used with the \texttt{svn\ merge} subcommand to merge
changes from a feature branch back into the feature
branch\textquotesingle s ancestor branch. Since Subversion 1.8 the
\texttt{svn\ merge} subcommand automatically detects this scenario and
performs the appropriate merge. See
\hyperref[svn.branchmerge.basicmerging.reintegrate]{Reintegrating a
Branch} for details.
\item[\protect\phantomsection\label{svn.ref.svn.sw.relocate}{}\texttt{-\/-relocate}]
\emph{Deprecated}. When used with the \texttt{svn\ switch} subcommand,
changes the location of the repository that your working copy
references. The preferred approach as of Subversion 1.7, however, is to
use the \texttt{svn\ relocate} subcommand. See
\hyperref[svn.ref.svn.c.relocate]{svn relocate} for more details and an
example.
\item[\protect\phantomsection\label{svn.ref.svn.sw.remove}{}\texttt{-\/-remove}]
Used with \texttt{svn\ changelist} to disassociate---rather than
associate (which is the default operation)---the target(s) from a
changelist.
\item[\protect\phantomsection\label{svn.ref.svn.sw.reverse_diff}{}\texttt{-\/-reverse-diff}]
Causes \texttt{svn\ patch} to interpret the input patch instructions in
reverse---treating added lines as removed ones and vice-versa.
\item[\protect\phantomsection\label{svn.ref.svn.sw.revision}{}\texttt{-\/-revision}
(\texttt{-r}) \textless REV\textgreater{}]
Specifies a revision (or range of revisions) on with which to operate.
You can provide revision numbers, keywords, or dates (in curly braces)
as arguments to the revision option. If you wish to offer a range of
revisions, you can provide two revisions separated by a colon. For
example:

\begin{verbatim}
$ svn log -r 1729
$ svn log -r 1729:HEAD
$ svn log -r 1729:1744
$ svn log -r {2001-12-04}:{2002-02-17}
$ svn log -r 1729:{2002-02-17}
\end{verbatim}

See \hyperref[svn.tour.revs.keywords]{Revision Keywords} for more
information.
\item[\protect\phantomsection\label{svn.ref.svn.sw.revprop}{}\texttt{-\/-revprop}]
Operates on a revision property instead of a property specific to a file
or directory. This option requires that you also pass a revision with
the \texttt{-\/-revision} (\texttt{-r}) option.
\item[\protect\phantomsection\label{svn.ref.svn.sw.search}{}\texttt{-\/-search}
\textless ARG\textgreater{}]
Filters log messages to show only those that match the search pattern
\textless ARG\textgreater. Log messages are displayed only if the
provided search pattern matches any of the author, date, log message
text (unless \texttt{-\/-quiet} is used), or, if the
\texttt{-\/-verbose} option is also provided, a changed path. If
multiple \texttt{-\/-search} options are provided, a log message is
shown if it matches any of the provided search patterns. If
\texttt{-\/-limit} is used, it restricts the number of log messages
searched, rather than restricting the output to a particular number of
matching log messages.

The search pattern (also called glob or shell wildcard pattern) can
contain regular characters and the following wildcard characters:

\begin{description}
\item[\texttt{?}]
Matches any single character.
\item[\texttt{*}]
Matches a sequence of arbitrary characters.
\item[\texttt{{[}ABC{]}}]
Matches any of the characters listed inside the brackets.
\end{description}
\item[\protect\phantomsection\label{svn.ref.svn.sw.search_and}{}\texttt{-\/-search-and}
\textless ARG\textgreater{}]
The option\textquotesingle s argument is combined with the pattern from
the previous \texttt{-\/-search} or \texttt{-\/-search-and} option on
the command line. Log message is shown only if it matches the combined
search pattern.
\item[\protect\phantomsection\label{svn.ref.svn.sw.set_depth}{}\texttt{-\/-set-depth}
\textless ARG\textgreater{}]
Sets the sticky depth on a directory in a working copy to one of
\texttt{exclude}, \texttt{empty}, \texttt{files}, \texttt{immediates},
or \texttt{infinity}. For detailed coverage of what these mean and how
to use this option, see \hyperref[svn.advanced.sparsedirs]{Sparse
Directories}.
\item[\protect\phantomsection\label{svn.ref.svn.sw.show_copies_as_adds}{}\texttt{-\/-show-copies-as-adds}]
Enables a special output mode for \texttt{svn\ diff} in which the
content difference for a file created via a copy operation appears as it
would for a brand new file (with each line therein appearing as an
addition to an empty file) rather than as a delta against the original
file from which the copy was created.
\item[\protect\phantomsection\label{svn.ref.svn.sw.show_inherited_props}{}\texttt{-\/-show-inherited-props}]
Causes \texttt{svn\ propget} and \texttt{svn\ proplist} to display the
versioned properties inherited by the target file or directory.
\item[\protect\phantomsection\label{svn.ref.svn.sw.show_revs}{}\texttt{-\/-show-revs}
\textless ARG\textgreater{}]
Used to make \texttt{svn\ mergeinfo} display certain classes of merge
tracking information. \textless ARG\textgreater{} may be either
\texttt{merged} or \texttt{eligible}, indicating a desire to see
revisions either already merged or eligible for future merge from the
specified source URL, respectively.
\item[\protect\phantomsection\label{svn.ref.svn.sw.show_updates}{}\texttt{-\/-show-updates}
(\texttt{-u})]
Causes the client to display information about which files in your
working copy are out of date. This doesn\textquotesingle t actually
update any of your files---it just shows you which files will be updated
if you then use \texttt{svn\ update}.
\item[\protect\phantomsection\label{svn.ref.svn.sw.stop_on_copy}{}\texttt{-\/-stop-on-copy}]
Causes a Subversion subcommand that traverses the history of a versioned
resource to stop harvesting that historical information when a
copy---that is, a location in history where that resource was copied
from another location in the repository---is encountered.
\item[\protect\phantomsection\label{svn.ref.svn.sw.strict}{}\texttt{-\/-strict}]
Causes Subversion to use strict semantics, a notion that is rather vague
unless talking about specific subcommands (namely,
\texttt{svn\ propget}).
\item[\protect\phantomsection\label{svn.ref.svn.sw.strip}{}\texttt{-\/-strip}
\textless NUM\textgreater{}]
Used by \texttt{svn\ patch} to ignore \textless NUM\textgreater{}
leading path components found on paths specified in the patch input
file.
\item[\protect\phantomsection\label{svn.ref.svn.sw.summarize}{}\texttt{-\/-summarize}]
Display only high-level summary notifications about the operation
instead of its detailed output.
\item[\protect\phantomsection\label{svn.ref.svn.sw.targets}{}\texttt{-\/-targets}
\textless FILENAME\textgreater{}]
Tells Subversion to read additional target paths for the operation from
\textless FILENAME\textgreater. \textless FILENAME\textgreater{} should
contain one path per line, with each path expected to use the same
encoding and formatting that it would if you had specified it directly
as an argument on the command line.
\item[\protect\phantomsection\label{svn.ref.svn.sw.trust_server_cert}{}\texttt{-\/-trust-server-cert}]
When used with \texttt{-\/-non-interactive}, instructs Subversion to
accept SSL server certificates issued by unknown certificate authorities
without first prompting the user. For security\textquotesingle s sake,
you should use this option only when the integrity of the remote server
and the network path between it and your client is known to be
trustworthy.

This option is accepted by all \texttt{svn} subcommands.
\item[\protect\phantomsection\label{svn.ref.svn.sw.use_merge_history}{}\texttt{-\/-use-merge-history}
(\texttt{-g})]
Uses or displays additional information from merge history.
\item[\protect\phantomsection\label{svn.ref.svn.sw.username}{}\texttt{-\/-username}
\textless NAME\textgreater{}]
Specifies the username to use when authenticating against a Subversion
server. If not provided, or if incorrect, Subversion will prompt you for
this information as needed.

This option is accepted by all \texttt{svn} subcommands.
\item[\protect\phantomsection\label{svn.ref.svn.sw.verbose}{}\texttt{-\/-verbose}
(\texttt{-v})]
Requests that the client print out as much information as it can while
running any subcommand. This may result in Subversion printing out
additional fields, detailed information about every file, or additional
information regarding its actions.
\item[\protect\phantomsection\label{svn.ref.svn.sw.version}{}\texttt{-\/-version}]
Prints the client version info. This information includes not only the
version number of the client, but also a listing of all repository
access modules that the client can use to access a Subversion
repository. With \texttt{-\/-quiet} (\texttt{-q}) it prints only the
version number in a compact form.
\item[\protect\phantomsection\label{svn.ref.svn.sw.with_all_revprops}{}\texttt{-\/-with-all-revprops}]
Used with the \texttt{-\/-xml} option to \texttt{svn\ log}, instructs
Subversion to retrieve and display all revision properties---the
standard ones used internally by Subversion as well as any user-defined
ones---in the log output.
\item[\protect\phantomsection\label{svn.ref.svn.sw.with_no_revprops}{}\texttt{-\/-with-no-revprops}]
Used with the \texttt{-\/-xml} option to \texttt{svn\ log}, instructs
Subversion to omit all revision properties---including the standard log
message, author, and revision datestamp---from the log output.
\item[\protect\phantomsection\label{svn.ref.svn.sw.with_revprop}{}\texttt{-\/-with-revprop}
\textless ARG\textgreater{}]
When used with any command that writes to the repository, sets the
revision property, using the \textless NAME=VALUE\textgreater{} format,
\textless NAME\textgreater{} to \textless VALUE\textgreater. When used
with \texttt{svn\ log} in \texttt{-\/-xml} mode, this displays the value
of \textless ARG\textgreater{} in the log output.
\item[\protect\phantomsection\label{svn.ref.svn.sw.xml}{}\texttt{-\/-xml}]
Prints output in XML format. XML schemas for the output (in RELAX NG
format) are maintained in the \texttt{subversion/svn/schema/} directory
of the Subversion source tree.
\end{description}

\section{svn add}\label{svn.ref.svn.c.add}

\emph{Add files, directories, or symbolic links.}

\textbf{Synopsis}

\texttt{svn\ add\ PATH...}

\subsection{Description}

Schedule files, directories, or symbolic links in your working copy for
addition to the repository. They will be uploaded and added to the
repository on your next commit. If you add something and change your
mind before committing, you can unschedule the addition using
\texttt{svn\ revert}.

\subsection{Options}

\begin{verbatim}
--auto-props
--depth ARG
--force
--no-auto-props
--no-ignore
--parents
--quiet (-q)
--targets FILENAME
\end{verbatim}

\subsection{Examples}

To add a file to your working copy:

\begin{verbatim}
$ svn add foo.c 
A         foo.c
\end{verbatim}

When adding a directory, the default behavior of \texttt{svn\ add} is to
recurse:

\begin{verbatim}
$ svn add testdir
A         testdir
A         testdir/a
A         testdir/b
A         testdir/c
A         testdir/d
\end{verbatim}

You can add a directory without adding its contents:

\begin{verbatim}
$ svn add --depth=empty otherdir
A         otherdir
\end{verbatim}

Attempts to schedule the addition of an item which is already versioned
will fail by default. This behavior foils the most common scenario under
which users attempt this: when trying to get to Subversion to
recursively examine a versioned directory and add any unversioned items
inside of it. To override the default behavior and force Subversion to
recurse into already-versioned directories, pass the \texttt{-\/-force}
option:

\begin{verbatim}
$ svn add versioned-dir
svn: warning: W150002: '/home/cmpilato/projects/subversion/site' is already un\
der version control
$ svn add versioned-dir --force
A         versioned-dir/foo.c
A         versioned-dir/somedir/bar.c
A  (bin)  versioned-dir/otherdir/docs/baz.doc
…
\end{verbatim}

\section{svn blame (praise, annotate, ann)}\label{svn.ref.svn.c.blame}

\emph{Show author and revision information inline for the specified
files or URLs.}

\textbf{Synopsis}

\texttt{svn\ blame\ TARGET{[}@REV{]}...}

\subsection{Description}

Show author and revision information inline for the specified files or
URLs. Each line of text is annotated at the beginning with the author
(username) and the revision number for the last change to that line.

\subsection{Options}

\begin{verbatim}
--extensions (-x) ARG
--force
--incremental
--revision (-r) REV
--use-merge-history (-g)
--verbose (-v)
--xml
\end{verbatim}

\subsection{Examples}

If you want to see blame-annotated source for \texttt{readme.txt} in
your test repository:

\begin{verbatim}
$ svn blame http://svn.red-bean.com/repos/test/readme.txt
     3      sally This is a README file.
     5      harry Don't bother reading it.  The boss is a knucklehead.
     3      sally 
…
\end{verbatim}

Now, just because \texttt{svn\ blame} says that Harry last modified
\texttt{readme.txt} in revision 5, understand that this subcommand is by
default very picky about what constitutes a change. Before clubbing
Harry over the head for what appears to be insubordination, first
consider that perhaps the change he made to the file might have been
only to its specific character content, not to its overall semantic
meaning. Perhaps his changes were the result of blindly running a
whitespace cleanup script on this file. You might need to examine the
specific differences and related log message to understand exactly what
Harry did to this file in revision 5.

\begin{verbatim}
$ svn log -c 5 http://svn.red-bean.com/repos/test/readme.txt
------------------------------------------------------------------------
r5 | harry | 2008-05-29 07:26:12 -0600 (Thu, 29 May 2008) | 1 line

Commit the results of 'double-space-after-period.sh'.

------------------------------------------------------------------------
$ svn diff -c 5 http://svn.red-bean.com/repos/test/readme.txt
Index: http://svn.red-bean.com/repos/test/readme.txt
===================================================================
--- http://svn.red-bean.com/repos/test/readme.txt   (revision 4)
+++ http://svn.red-bean.com/repos/test/readme.txt   (revision 5)
@@ -1,5 +1,5 @@
 This is a README file.
-Don't bother reading it. The boss is a knucklehead.
+Don't bother reading it.  The boss is a knucklehead.
  
 INSTRUCTIONS
 ============
$
\end{verbatim}

Sure enough, Harry only changed the whitespace in that line.
Fortunately, the \texttt{-\/-extensions} (\texttt{-x}) option can help
you better determine the last time that a \emph{meaningful} change was
made to a given line of text. For example, here\textquotesingle s how
you can see the annotation information while disregarding mere
whitespace changes:

\begin{verbatim}
$ svn blame -x -b http://svn.red-bean.com/repos/test/readme.txt
     3      sally This is a README file.
     4       jess Don't bother reading it.  The boss is a knucklehead.
     3      sally 
…
\end{verbatim}

If you use the \texttt{-\/-xml} option, you can get XML output
describing the blame annotations, but not the contents of the lines
themselves:

\begin{verbatim}
$ svn blame --xml http://svn.red-bean.com/repos/test/readme.txt
<?xml version="1.0"?>
<blame>
<target
   path="readme.txt">
<entry
   line-number="1">
<commit
   revision="3">
<author>sally</author>
<date>2008-05-25T19:12:31.428953Z</date>
</commit>
</entry>
<entry
   line-number="2">
<commit
   revision="5">
<author>harry</author>
<date>2008-05-29T13:26:12.293121Z</date>
</commit>
</entry>
<entry
   line-number="3">
…
</entry>
</target>
</blame>
$
\end{verbatim}

\section{svn cat}\label{svn.ref.svn.c.cat}

\emph{Output the contents of the specified files or URLs.}

\textbf{Synopsis}

\texttt{svn\ cat\ TARGET{[}@REV{]}...}

\subsection{Description}

Output the contents of the specified files or URLs. For listing the
contents of directories, see \texttt{svn\ list} later in this chapter.

\subsection{Options}

\begin{verbatim}
--revision (-r) REV
\end{verbatim}

\subsection{Examples}

If you want to view \texttt{readme.txt} in your repository without
checking it out:

\begin{verbatim}
$ svn cat http://svn.red-bean.com/repos/test/readme.txt
This is a README file.
Don't bother reading it.  The boss is a knucklehead.
 
INSTRUCTIONS
============

Step 1:  Do this.

Step 2:  Do that.
$
\end{verbatim}

You can view specific versions of files, too.

\begin{verbatim}
$ svn cat -r 3 http://svn.red-bean.com/repos/test/readme.txt
This is a README file.
 
INSTRUCTIONS
============

Step 1:  Do this.

Step 2:  Do that.
$
\end{verbatim}

You might develop a reflex action of using \texttt{svn\ cat} to view
your working file contents. But keep in mind that the default peg
revision for \texttt{svn\ cat} when used on a working copy file target
is \texttt{BASE}, the unmodified base revision of that file.
Don\textquotesingle t be surprised when a simple
\texttt{svn\ cat\ /path/to/file} invocation fails to display your local
modifications to that file!

If your working copy is out of date (or you have local modifications)
and you want to see the \texttt{HEAD} revision of a file in your working
copy, use the \texttt{-\/-revision} (\texttt{-r}) option:
\texttt{svn\ cat\ -r\ HEAD\ FILENAME}

\section{svn changelist (cl)}\label{svn.ref.svn.c.changelist}

\emph{Associate (or deassociate) local paths with a changelist.}

\textbf{Synopsis}

\texttt{changelist\ CLNAME\ TARGET...}

\texttt{changelist\ -\/-remove\ TARGET...}

\subsection{Description}

Used for dividing files in a working copy into a changelist (logical
named grouping) in order to allow users to easily work on multiple file
collections within a single working copy.

\subsection{Options}

\begin{verbatim}
--changelist (--cl) ARG
--depth ARG
--quiet (-q)
--recursive (-R)
--remove
--targets FILENAME
\end{verbatim}

\subsection{Example}

Edit three files, add them to a changelist, then commit only files in
that changelist:

\begin{verbatim}
$ svn changelist issue1729 foo.c bar.c baz.c
A [issue1729] foo.c
A [issue1729] bar.c
A [issue1729] baz.c
$ svn status
A       someotherfile.c
A       test/sometest.c

--- Changelist 'issue1729':
A       foo.c
A       bar.c
A       baz.c
$ svn commit --changelist issue1729 -m "Fixing Issue 1729."
Adding         bar.c
Adding         baz.c
Adding         foo.c
Transmitting file data ...
Committed revision 2.
$ svn status
A       someotherfile.c
A       test/sometest.c
$
\end{verbatim}

Note that in the previous example, only the files in changelist
\texttt{issue1729} were committed.

\section{svn checkout (co)}\label{svn.ref.svn.c.checkout}

\emph{Check out a working copy from a repository.}

\textbf{Synopsis}

\texttt{svn\ checkout\ URL{[}@REV{]}...\ {[}PATH{]}}

\subsection{Description}

Check out a working copy from a repository. If
\textless PATH\textgreater{} is omitted, the basename of the URL will be
used as the destination. If multiple URLs are given, each will be
checked out into a subdirectory of \textless PATH\textgreater, with the
name of the subdirectory being the basename of the URL.

\subsection{Options}

\begin{verbatim}
--depth ARG
--force
--ignore-externals
--quiet (-q)
--revision (-r) REV
\end{verbatim}

\subsection{Examples}

Check out a working copy into a directory called \texttt{mine}:

\begin{verbatim}
$ svn checkout file:///var/svn/repos/test mine
A    mine/a
A    mine/b
A    mine/c
A    mine/d
Checked out revision 20.
$ ls
mine
$
\end{verbatim}

Check out two different directories into two separate working copies:

\begin{verbatim}
$ svn checkout file:///var/svn/repos/test \
               file:///var/svn/repos/quiz
A    test/a
A    test/b
A    test/c
A    test/d
Checked out revision 20.
A    quiz/l
A    quiz/m
Checked out revision 13.
$ ls
quiz  test
$
\end{verbatim}

Check out two different directories into two separate working copies,
but place both into a directory called \texttt{working-copies}:

\begin{verbatim}
$ svn checkout file:///var/svn/repos/test \
               file:///var/svn/repos/quiz \
               working-copies
A    working-copies/test/a
A    working-copies/test/b
A    working-copies/test/c
A    working-copies/test/d
Checked out revision 20.
A    working-copies/quiz/l
A    working-copies/quiz/m
Checked out revision 13.
$ ls
working-copies
\end{verbatim}

If you interrupt a checkout (or something else interrupts your checkout,
such as loss of connectivity, etc.), you can restart it either by
issuing the identical checkout command again or by updating the
incomplete working copy:

\begin{verbatim}
$ svn checkout file:///var/svn/repos/test mine
A    mine/a
A    mine/b
^C
svn: E200015: Caught signal
$ svn checkout file:///var/svn/repos/test mine
A    mine/c
^C
svn: E200015: Caught signal
$ svn update mine
Updating 'mine':
A    mine/d
Updated to revision 20.
$
\end{verbatim}

If you wish to check out some revision other than the most recent one,
you can do so by providing the \texttt{-\/-revision} (\texttt{-r})
option to the \texttt{svn\ checkout} command:

\begin{verbatim}
$ svn checkout -r 2 file:///var/svn/repos/test mine
A    mine/a
Checked out revision 2.
$
\end{verbatim}

Prior to version 1.7, Subversion would complain by default if you try to
check out a directory atop an existing directory which contains files or
subdirectories that the checkout itself would have created. Subversion
1.7 handles this situation differently, allowing the checkout to proceed
but marking any obstructing objects as tree conflicts. Use the
\texttt{-\/-force} option to override this safeguard. When you check out
with the \texttt{-\/-force} option, any unversioned file in the checkout
target tree which ordinarily would obstruct the checkout will still
become versioned, but Subversion will preserve its contents as-is. If
those contents differ from the repository file at that path (which was
downloaded as part of the checkout), the file will appear to have local
modifications---the changes required to transform the versioned file you
checked out into the unversioned file you had before checking out---when
the checkout completes.

\begin{verbatim}
$ mkdir project
$ mkdir project/lib
$ touch project/lib/file.c
$ svn checkout file:///var/svn/repos/project/trunk project --force
E    project/lib
A    project/lib/subdir
E    project/lib/file.c
A    project/lib/anotherfile.c
A    project/include/header.h
Checked out revision 21.
$ svn status wc
M       project/lib/file.c
$ svn diff wc
Index: project/lib/file.c
===================================================================
--- project/lib/file.c  (revision 1)
+++ project/lib/file.c  (working copy)
@@ -3 +0,0 @@
-/* file.c: Code for acting file-ishly. */
-#include <stdio.h>
-/* Not feeling particularly creative today. */

$
\end{verbatim}

As in any other working copy, you have the choices typically available:
reverting some or all of those local ``modifications'', committing them,
or continuing to modify your working copy.

This feature is especially useful for performing in-place imports of
unversioned directory trees. By first importing the tree into the
repository, and then checking out new repository location atop the
unversioned tree with the \texttt{-\/-force} option, you effectively
transform the unversioned tree into a working copy.

\begin{verbatim}
$ svn mkdir -m "Create newproject project root." \
      file://var/svn/repos/newproject
$ svn import -m "Import initial newproject codebase." newproject \
      file://var/svn/repos/newproject/trunk
Adding         newproject/include
Adding         newproject/include/newproject.h
Adding         newproject/lib
Adding         newproject/lib/helpers.c
Adding         newproject/lib/base.c
Adding         newproject/notes
Adding         newproject/notes/README

Committed revision 22.
$ svn checkout file://`pwd`/repos-1.6/newproject/trunk newproject --force
E    newproject/include
E    newproject/include/newproject.h
E    newproject/lib
E    newproject/lib/helpers.c
E    newproject/lib/base.c
E    newproject/notes
E    newproject/notes/README
Checked out revision 2.
$ svn status newproject
$
\end{verbatim}

\section{svn cleanup}\label{svn.ref.svn.c.cleanup}

\emph{Recursively clean up the working copy}

\textbf{Synopsis}

\texttt{svn\ cleanup\ {[}PATH...{]}}

\subsection{Description}

Recursively clean up the working copy, removing working copy locks and
resuming unfinished operations. If you ever get a
\texttt{working\ copy\ locked} error, run this command to remove stale
locks and get your working copy into a usable state again.

If, for some reason, an \texttt{svn\ update} fails due to a problem
running an external diff program (e.g., user input or network failure),
pass the \texttt{-\/-diff3-cmd} to allow the cleanup process to complete
any required merging using your external diff program. You can also
specify any configuration directory with the \texttt{-\/-config-dir}
option, but you should need these options extremely infrequently.

\subsection{Options}

\begin{verbatim}
--diff3-cmd CMD
\end{verbatim}

\subsection{Examples}

Well, there\textquotesingle s not much to the examples here, as
\texttt{svn\ cleanup} generates no output. If you pass no
\textless PATH\textgreater, then ``\texttt{.}'' is used:

\begin{verbatim}
$ svn cleanup
$ svn cleanup /var/svn/working-copy
\end{verbatim}

\section{svn commit (ci)}\label{svn.ref.svn.c.commit}

\emph{Send changes from your working copy to the repository.}

\textbf{Synopsis}

\texttt{svn\ commit\ {[}PATH...{]}}

\subsection{Description}

Send changes from your working copy to the repository. If you do not
supply a log message with your commit by using either the
\texttt{-\/-file} (\texttt{-F}) or \texttt{-\/-message} (\texttt{-m})
option, \texttt{svn} will launch your editor for you to compose a commit
message. See the \texttt{editor-cmd} list entry in
\hyperref[svn.advanced.confarea.opts.config]{General configuration}.

\texttt{svn\ commit} will send any lock tokens that it finds and will
release locks on all \textless PATH\textgreater s committed
(recursively) unless \texttt{-\/-no-unlock} is passed.

If you begin a commit and Subversion launches your editor to compose the
commit message, you can still abort without committing your changes. If
you want to cancel your commit, just quit your editor without saving
your commit message and Subversion will prompt you to either abort the
commit, continue with no message, or edit the message again.

\subsection{Options}

\begin{verbatim}
--changelist (--cl) ARG
--depth ARG
--editor-cmd CMD
--encoding ENC
--file (-F) FILENAME
--force-log
--keep-changelists
--message (-m) MESSAGE
--no-unlock
--quiet (-q)
--targets FILENAME
--with-revprop ARG
\end{verbatim}

\subsection{Examples}

Commit a simple modification to a file with the commit message on the
command line and an implicit target of your current directory
(``\texttt{.}''):

\begin{verbatim}
$ svn commit -m "added howto section."
Sending        a
Transmitting file data .
Committed revision 3.
\end{verbatim}

Commit a modification to the file \texttt{foo.c} (explicitly specified
on the command line) with the commit message in a file named
\texttt{msg}:

\begin{verbatim}
$ svn commit -F msg foo.c
Sending        foo.c
Transmitting file data .
Committed revision 5.
\end{verbatim}

If you want to use a file that\textquotesingle s under version control
for your commit message with \texttt{-\/-file} (\texttt{-F}), you need
to pass the \texttt{-\/-force-log} option:

\begin{verbatim}
$ svn commit -F file_under_vc.txt foo.c
svn: E205004: Log message file is a versioned file; use '--force-log' to override

$ svn commit --force-log -F file_under_vc.txt foo.c
Sending        foo.c
Transmitting file data .
Committed revision 6.
\end{verbatim}

To commit a file scheduled for deletion:

\begin{verbatim}
$ svn commit -m "removed file 'c'."
Deleting       c

Committed revision 7.
\end{verbatim}

\section{svn copy (cp)}\label{svn.ref.svn.c.copy}

\emph{Copy a file or directory in a working copy or in the repository.}

\textbf{Synopsis}

\texttt{svn\ copy\ SRC{[}@REV{]}...\ DST}

\subsection{Description}

Copy one or more files in a working copy or in the repository.
\textless SRC\textgreater{} and \textless DST\textgreater{} can each be
either a working copy (WC) path or URL. When copying multiple sources,
add the copies as immediate children of \textless DST\textgreater{}
(which, of course, must be an existing directory).

\begin{description}
\item[WC → WC]
Copy and schedule an item for addition (with history).
\item[WC → URL]
Immediately commit a copy of WC to URL.
\item[URL → WC]
Check out URL into WC and schedule it for addition.
\item[URL → URL]
Complete server-side copy. This is usually used to branch and tag.
\end{description}

If no peg revision (i.e., \texttt{@REV}) is supplied, by default the
\texttt{BASE} revision will be used for files copied from the working
copy, while the \texttt{HEAD} revision will be used for files copied
from a URL.

You can only copy files within a single repository. Subversion does not
support cross-repository copying.

\subsection{Options}

\begin{verbatim}
--editor-cmd CMD
--encoding ENC
--file (-F) FILENAME
--force-log
--ignore-externals
--message (-m) MESSAGE
--parents
--quiet (-q)
--revision (-r) REV
--with-revprop ARG
\end{verbatim}

\subsection{Examples}

Copy an item within your working copy (this schedules the copy---nothing
goes into the repository until you commit):

\begin{verbatim}
$ svn copy foo.txt bar.txt
A         bar.txt
$ svn status
A  +    bar.txt
\end{verbatim}

Copy several files in a working copy into a directory:

\begin{verbatim}
$ svn copy bat.c baz.c qux.c src
A         src/bat.c
A         src/baz.c
A         src/qux.c
\end{verbatim}

Copy revision 8 of \texttt{bat.c} into your working copy under a
different name:

\begin{verbatim}
$ svn copy -r 8 bat.c ya-old-bat.c
A         ya-old-bat.c
\end{verbatim}

Copy an item in your working copy to a URL in the repository (this is an
immediate commit, so you must supply a commit message):

\begin{verbatim}
$ svn copy near.txt file:///var/svn/repos/test/far-away.txt -m "Remote copy."

Committed revision 8.
\end{verbatim}

Copy an item from the repository to your working copy (this just
schedules the copy---nothing goes into the repository until you commit):

\begin{verbatim}
$ svn copy file:///var/svn/repos/test/far-away -r 6 near-here
A         near-here
\end{verbatim}

This is the recommended way to resurrect a dead file in your repository!

And finally, copy between two URLs:

\begin{verbatim}
$ svn copy file:///var/svn/repos/test/far-away \
           file:///var/svn/repos/test/over-there -m "remote copy."

Committed revision 9.
\end{verbatim}

\begin{verbatim}
$ svn copy file:///var/svn/repos/test/trunk \
           file:///var/svn/repos/test/tags/0.6.32-prerelease -m "tag tree"

Committed revision 12.
\end{verbatim}

This is the easiest way to ``tag'' a revision in your repository---just
\texttt{svn\ copy} that revision (usually \texttt{HEAD}) into your
\texttt{tags} directory.

And don\textquotesingle t worry if you forgot to tag---you can always
specify an older revision and tag anytime:

\begin{verbatim}
$ svn copy -r 11 file:///var/svn/repos/test/trunk \
           file:///var/svn/repos/test/tags/0.6.32-prerelease \
           -m "Forgot to tag at rev 11"

Committed revision 13.
\end{verbatim}

\section{svn delete (del, remove, rm)}\label{svn.ref.svn.c.delete}

\emph{Delete an item from a working copy or the repository.}

\textbf{Synopsis}

\texttt{svn\ delete\ PATH...}

\texttt{svn\ delete\ URL...}

\subsection{Description}

Items specified by \textless PATH\textgreater{} are scheduled for
deletion upon the next commit. Files (and directories that have not been
committed) are immediately removed from the working copy unless the
\texttt{-\/-keep-local} option is given. The command will not remove any
unversioned or modified items; use the \texttt{-\/-force} option to
override this behavior. A directory that contains unversioned or
modified items will not be deleted, unless the \texttt{-\/-force} is
used.

Items specified by URL are deleted from the repository via an immediate
commit. Multiple URLs are committed atomically.

\subsection{Options}

\begin{verbatim}
--editor-cmd CMD
--encoding ENC
--file (-F) FILENAME
--force
--force-log
--keep-local
--message (-m) MESSAGE
--quiet (-q)
--targets FILENAME
--with-revprop ARG
\end{verbatim}

\subsection{Examples}

Using \texttt{svn} to delete a file from your working copy deletes your
local copy of the file, but it merely schedules the file to be deleted
from the repository. When you commit, the file is deleted in the
repository.

\begin{verbatim}
$ svn delete myfile
D         myfile

$ svn commit -m "Deleted file 'myfile'."
Deleting       myfile
Transmitting file data .
Committed revision 14.
\end{verbatim}

Deleting a URL, however, is immediate, so you have to supply a log
message:

\begin{verbatim}
$ svn delete -m "Deleting file 'yourfile'" \
             file:///var/svn/repos/test/yourfile

Committed revision 15.
\end{verbatim}

Here\textquotesingle s an example of how to force deletion of a file
that has local modifications:

\begin{verbatim}
$ svn delete over-there 
svn: E195006: Use --force to override this restriction (local modifications m\
ay be lost)
svn: E195006: '/home/sally/project/over-there' has local modifications -- com\
mit or revert them first
$ svn delete --force over-there 
D         over-there
$
\end{verbatim}

Use the \texttt{-\/-keep-local} option to override the default
\texttt{svn\ delete} behavior of also removing the target file that was
scheduled for versioned deletion. This is helpful when you realize that
you\textquotesingle ve accidentally committed the addition of a file
that you need to keep around in your working copy, but which
shouldn\textquotesingle t have been added to version control.

\begin{verbatim}
$ svn delete --keep-local conf/program.conf
D         conf/program.conf

$ svn commit -m "Remove accidentally-added configuration file."
Deleting       conf/program.conf
Transmitting file data .
Committed revision 21.
$ svn status
?       conf/program.conf
$
\end{verbatim}

The behavior of the \texttt{-\/-keep-local} option does not propagate to
other working copies which contain the items you\textquotesingle ve
scheduled for deletion. If you commit the deletion of those items they
will remain in your working copy, but they will be deleted from other
working copies which contain them when those working copies are then
updated.

\section{svn diff (di)}\label{svn.ref.svn.c.diff}

\emph{This displays the differences between two revisions or paths.}

\textbf{Synopsis}

\texttt{diff\ {[}-c\ M\ \textbar{}\ -r\ N{[}:M{]}{]}\ {[}TARGET{[}@REV{]}...{]}}

\texttt{diff\ {[}-r\ N{[}:M{]}{]}\ -\/-old=OLD-TGT{[}@OLDREV{]}\ {[}-\/-new=NEW-TGT{[}@NEWREV{]}{]}\ {[}PATH...{]}}

\texttt{diff\ OLD-URL{[}@OLDREV{]}\ NEW-URL{[}@NEWREV{]}}

\subsection{Description}

Display the differences between two paths. You can use
\texttt{svn\ diff} in the following ways:

\begin{itemize}
\item
  Use just \texttt{svn\ diff} to display local modifications in a
  working copy.
\item
  Display the changes made to \textless TARGET\textgreater s as they are
  seen in \textless REV\textgreater{} between two revisions.
  \textless TARGET\textgreater s may be all working copy paths or all
  \textless URL\textgreater s. If \textless TARGET\textgreater s are
  working copy paths, \textless N\textgreater{} defaults to
  \texttt{BASE} and \textless M\textgreater{} to the working copy; if
  \textless TARGET\textgreater s are \textless URL\textgreater s,
  \textless N\textgreater{} must be specified and
  \textless M\textgreater{} defaults to \texttt{HEAD}. The
  \texttt{-c\ M} option is equivalent to \texttt{-r\ N:M} where
  \texttt{N\ =\ M-1}. Using \texttt{-c\ -M} does the reverse:
  \texttt{-r\ M:N} where \texttt{N\ =\ M-1}.
\item
  Display the differences between \textless OLD-TGT\textgreater{} as it
  was seen in \textless OLDREV\textgreater{} and
  \textless NEW-TGT\textgreater{} as it was seen in
  \textless NEWREV\textgreater. \textless PATH\textgreater s, if given,
  are relative to \textless OLD-TGT\textgreater{} and
  \textless NEW-TGT\textgreater{} and restrict the output to differences
  for those paths. \textless OLD-TGT\textgreater{} and
  \textless NEW-TGT\textgreater{} may be working copy paths or
  \textless URL{[}@REV{]}\textgreater. \textless NEW-TGT\textgreater{}
  defaults to \textless OLD-TGT\textgreater{} if not specified.
  \texttt{-r\ N} makes \textless OLDREV\textgreater{} default to
  \texttt{N}; \texttt{-r\ N:M} makes \textless OLDREV\textgreater{}
  default to \textless N\textgreater{} and
  \textless NEWREV\textgreater{} default to \textless M\textgreater.
\end{itemize}

\texttt{svn\ diff\ OLD-URL{[}@OLDREV{]}\ NEW-URL{[}@NEWREV{]}} is
shorthand for
\texttt{svn\ diff\ -\/-old=OLD-URL{[}@OLDREV{]}\ -\/-new=NEW-URL{[}@NEWREV{]}}.

\texttt{svn\ diff\ -r\ N:M\ URL} is shorthand for
\texttt{svn\ diff\ -r\ N:M\ -\/-old=URL\ -\/-new=URL}.

\texttt{svn\ diff\ {[}-r\ N{[}:M{]}{]}\ URL1{[}@N{]}\ URL2{[}@M{]}} is
shorthand for
\texttt{svn\ diff\ {[}-r\ N{[}:M{]}{]}\ -\/-old=URL1\ -\/-new=URL2}.

If \textless TARGET\textgreater{} is a URL, then revisions
\textless N\textgreater{} and \textless M\textgreater{} can be given
either via the \texttt{-\/-revision} (\texttt{-r}) option or by using
the ``@'' notation as described earlier.

If \textless TARGET\textgreater{} is a working copy path, the default
behavior (when no \texttt{-\/-revision} (\texttt{-r}) option is
provided) is to display the differences between the base and working
copies of \textless TARGET\textgreater. If a \texttt{-\/-revision}
(\texttt{-r}) option is specified in this scenario, though, it means:

\begin{description}
\item[\texttt{-\/-revision\ N:M}]
The server compares \textless TARGET@N\textgreater{} and
\textless TARGET@M\textgreater.
\item[\texttt{-\/-revision\ N}]
The client compares \textless TARGET@N\textgreater{} against the working
copy.
\end{description}

If the alternate syntax is used, the server compares
\textless URL1\textgreater{} and \textless URL2\textgreater{} at
revisions \textless N\textgreater{} and \textless M\textgreater,
respectively. If either \textless N\textgreater{} or
\textless M\textgreater{} is omitted, a value of \texttt{HEAD} is
assumed.

By default, \texttt{svn\ diff} ignores the ancestry of files and merely
compares the contents of the two files being compared. If you use
\texttt{-\/-notice-ancestry}, the ancestry of the paths in question will
be taken into consideration when comparing revisions (i.e., if you run
\texttt{svn\ diff} on two files with identical contents but different
ancestry, you will see the entire contents of the file as having been
removed and added again).

\subsection{Options}

\begin{verbatim}
--change (-c) ARG
--changelist (--cl) ARG
--depth ARG
--diff-cmd CMD
--extensions (-x) ARG
--force
--git
--ignore-properties
--internal-diff
--new ARG
--no-diff-added
--no-diff-deleted
--notice-ancestry
--old ARG
--patch-compatible
--properties-only
--revision (-r) REV
--show-copies-as-adds
--summarize
--xml
\end{verbatim}

\subsection{Examples}

Compare \texttt{BASE} and your working copy (one of the most popular
uses of \texttt{svn\ diff}):

\begin{verbatim}
$ svn diff COMMITTERS 
Index: COMMITTERS
===================================================================
--- COMMITTERS  (revision 4404)
+++ COMMITTERS  (working copy)
…
\end{verbatim}

See what changed in the file \texttt{COMMITTERS} revision 9115:

\begin{verbatim}
$ svn diff -c 9115 COMMITTERS 
Index: COMMITTERS
===================================================================
--- COMMITTERS  (revision 3900)
+++ COMMITTERS  (working copy)
…
\end{verbatim}

See how your working copy\textquotesingle s modifications compare
against an older revision:

\begin{verbatim}
$ svn diff -r 3900 COMMITTERS 
Index: COMMITTERS
===================================================================
--- COMMITTERS  (revision 3900)
+++ COMMITTERS  (working copy)
…
\end{verbatim}

Compare revision 3000 to revision 3500 using ``@'' syntax:

\begin{verbatim}
$ svn diff http://svn.collab.net/repos/svn/trunk/COMMITTERS@3000 \
           http://svn.collab.net/repos/svn/trunk/COMMITTERS@3500
Index: COMMITTERS
===================================================================
--- COMMITTERS  (revision 3000)
+++ COMMITTERS  (revision 3500)
…
\end{verbatim}

Compare revision 3000 to revision 3500 using range notation (pass only
the one URL in this case):

\begin{verbatim}
$ svn diff -r 3000:3500 http://svn.collab.net/repos/svn/trunk/COMMITTERS
Index: COMMITTERS
===================================================================
--- COMMITTERS  (revision 3000)
+++ COMMITTERS  (revision 3500)
…
\end{verbatim}

Compare revision 3000 to revision 3500 of all the files in
\texttt{trunk} using range notation:

\begin{verbatim}
$ svn diff -r 3000:3500 http://svn.collab.net/repos/svn/trunk
\end{verbatim}

Compare revision 3000 to revision 3500 of only three files in
\texttt{trunk} using range notation:

\begin{verbatim}
$ svn diff -r 3000:3500 --old http://svn.collab.net/repos/svn/trunk \
       COMMITTERS README HACKING
\end{verbatim}

If you have a working copy, you can obtain the differences without
typing in the long URLs:

\begin{verbatim}
$ svn diff -r 3000:3500 COMMITTERS 
Index: COMMITTERS
===================================================================
--- COMMITTERS  (revision 3000)
+++ COMMITTERS  (revision 3500)
…
\end{verbatim}

Use \texttt{-\/-diff-cmd} \textless CMD\textgreater{}
\texttt{-\/-extensions} (\texttt{-x}) to pass arguments directly to the
external diff program:

\begin{verbatim}
$ svn diff --diff-cmd /usr/bin/diff -x "-i -b" COMMITTERS 
Index: COMMITTERS
===================================================================
0a1,2
> This is a test
> 
$
\end{verbatim}

Lastly, you can use the \texttt{-\/-xml} option along with the
\texttt{-\/-summarize} option to view XML describing the changes that
occurred between revisions, but not the contents of the diff itself:

\begin{verbatim}
$ svn diff --summarize --xml http://svn.red-bean.com/repos/test@r2 \
           http://svn.red-bean.com/repos/test
<?xml version="1.0"?>
<diff>
<paths>
<path
   props="none"
   kind="file"
   item="modified">http://svn.red-bean.com/repos/test/sandwich.txt</path>
<path
   props="none"
   kind="file"
   item="deleted">http://svn.red-bean.com/repos/test/burrito.txt</path>
<path
   props="none"
   kind="dir"
   item="added">http://svn.red-bean.com/repos/test/snacks</path>
</paths>
</diff>
\end{verbatim}

\section{svn export}\label{svn.ref.svn.c.export}

\emph{Export a clean directory tree.}

\textbf{Synopsis}

\texttt{svn\ export\ {[}-r\ REV{]}\ URL{[}@PEGREV{]}\ {[}PATH{]}}

\texttt{svn\ export\ {[}-r\ REV{]}\ PATH1{[}@PEGREV{]}\ {[}PATH2{]}}

\subsection{Description}

The first form exports a clean directory tree from the repository
specified by \textless URL\textgreater---at revision
\textless REV\textgreater{} if it is given; otherwise, at \texttt{HEAD},
into \textless PATH\textgreater. If \textless PATH\textgreater{} is
omitted, the last component of the \textless URL\textgreater{} is used
for the local directory name.

The second form exports a clean directory tree from the working copy
specified by \textless PATH1\textgreater{} into
\textless PATH2\textgreater. All local changes will be preserved, but
files not under version control will not be copied.

\subsection{Options}

\begin{verbatim}
--depth ARG
--force
--ignore-externals
--ignore-keywords
--native-eol ARG
--quiet (-q)
--revision (-r) REV
\end{verbatim}

\subsection{Examples}

Export from your working copy (doesn\textquotesingle t print every file
and directory):

\begin{verbatim}
$ svn export a-wc my-export
Export complete.
\end{verbatim}

Export directly from the repository (prints every file and directory):

\begin{verbatim}
$ svn export file:///var/svn/repos my-export
A    my-export/test
A    my-export/quiz
…
Exported revision 15.
\end{verbatim}

When rolling operating-system-specific release packages, it can be
useful to export a tree that uses a specific EOL character for line
endings. The \texttt{-\/-native-eol} option will do this, but it affects
only files that have \texttt{svn:eol-style\ =\ native} properties
attached to them. For example, to export a tree with all CRLF line
endings (possibly for a Windows \texttt{.zip} file distribution):

\begin{verbatim}
$ svn export file:///var/svn/repos my-export --native-eol CRLF
A    my-export/test
A    my-export/quiz
…
Exported revision 15.
\end{verbatim}

You can specify \texttt{LR}, \texttt{CR}, or \texttt{CRLF} as a
line-ending type with the \texttt{-\/-native-eol} option.

\section{svn help (h, ?)}\label{svn.ref.svn.c.help}

\emph{Help!}

\textbf{Synopsis}

\texttt{svn\ help\ {[}SUBCOMMAND...{]}}

\subsection{Description}

This is your best friend when you\textquotesingle re using Subversion
and this book isn\textquotesingle t within reach!

\subsection{Options}

None

\section{svn import}\label{svn.ref.svn.c.import}

\emph{Commit an unversioned file or tree into the repository.}

\textbf{Synopsis}

\texttt{svn\ import\ {[}PATH{]}\ URL}

\subsection{Description}

Recursively commit a copy of \textless PATH\textgreater{} to
\textless URL\textgreater. If \textless PATH\textgreater{} is omitted,
``\texttt{.}'' is assumed. Parent directories are created in the
repository as necessary. Unversionable items such as device files and
pipes are ignored even if \texttt{-\/-force} is specified.

\subsection{Options}

\begin{verbatim}
--auto-props
--depth ARG
--editor-cmd CMD
--encoding ENC
--file (-F) FILENAME
--force
--force-log
--message (-m) MESSAGE
--no-auto-props
--no-ignore
--quiet (-q)
--with-revprop ARG
\end{verbatim}

\subsection{Examples}

This imports the local directory \texttt{myproj} into
\texttt{trunk/misc} in your repository. The directory
\texttt{trunk/misc} need not exist before you import into
it---\texttt{svn\ import} will recursively create directories for you.

\begin{verbatim}
$ svn import -m "New import" myproj \
             http://svn.red-bean.com/repos/trunk/misc
Adding         myproj/sample.txt
…
Transmitting file data .........
Committed revision 16.
\end{verbatim}

Be aware that this will \emph{not} create a directory named
\texttt{myproj} in the repository. If that\textquotesingle s what you
want, simply add \texttt{myproj} to the end of the URL:

\begin{verbatim}
$ svn import -m "New import" myproj \
            http://svn.red-bean.com/repos/trunk/misc/myproj
Adding         myproj/sample.txt
…
Transmitting file data .........
Committed revision 16.
\end{verbatim}

After importing data, note that the original tree is \emph{not} under
version control. To start working, you still need to
\texttt{svn\ checkout} a fresh working copy of the tree.

\section{svn info}\label{svn.ref.svn.c.info}

\emph{Display information about a local or remote item.}

\textbf{Synopsis}

\texttt{svn\ info\ {[}TARGET{[}@REV{]}...{]}}

\subsection{Description}

Print information about the working copy paths or URLs specified. The
information displayed for each path may include (as pertinent to the
object at that path):

\begin{itemize}
\item
  information about the repository in which the object is versioned
\item
  the most recent commit made to the specified version of the object
\item
  any user-level locks held on the object
\item
  local scheduling information (added, deleted, copied, etc.)
\item
  local conflict information
\end{itemize}

\subsection{Options}

\begin{verbatim}
--changelist (--cl) ARG
--depth ARG
--incremental
--recursive (-R)
--revision (-r) REV
--targets FILENAME
--xml
\end{verbatim}

\subsection{Examples}

\texttt{svn\ info} will show you all the useful information that it has
for items in your working copy. It will show information for files:

\begin{verbatim}
$ svn info foo.c
Path: foo.c
Name: foo.c
Working Copy Root Path: /home/sally/projects/test
URL: http://svn.red-bean.com/repos/test/foo.c
Relative URL: ^/foo.c
Repository Root: http://svn.red-bean.com/repos/test
Repository UUID: 5e7d134a-54fb-0310-bd04-b611643e5c25
Revision: 4417
Node Kind: file
Schedule: normal
Last Changed Author: sally
Last Changed Rev: 20
Last Changed Date: 2003-01-13 16:43:13 -0600 (Mon, 13 Jan 2003)
Text Last Updated: 2003-01-16 21:18:16 -0600 (Thu, 16 Jan 2003)
Properties Last Updated: 2003-01-13 21:50:19 -0600 (Mon, 13 Jan 2003)
Checksum: d6aeb60b0662ccceb6bce4bac344cb66
\end{verbatim}

It will also show information for directories:

\begin{verbatim}
$ svn info vendors
Path: vendors
Working Copy Root Path: /home/sally/projects/test
URL: http://svn.red-bean.com/repos/test/vendors
Relative URL: ^/vendors
Repository Root: http://svn.red-bean.com/repos/test
Repository UUID: 5e7d134a-54fb-0310-bd04-b611643e5c25
Revision: 19
Node Kind: directory
Schedule: normal
Last Changed Author: harry
Last Changed Rev: 19
Last Changed Date: 2003-01-16 23:21:19 -0600 (Thu, 16 Jan 2003)
Properties Last Updated: 2003-01-16 23:39:02 -0600 (Thu, 16 Jan 2003)
\end{verbatim}

\texttt{svn\ info} also acts on URLs (also note that the file
\texttt{readme.doc} in this example is locked, so lock information is
also provided):

\begin{verbatim}
$ svn info http://svn.red-bean.com/repos/test/readme.doc
Path: readme.doc
Name: readme.doc
URL: http://svn.red-bean.com/repos/test/readme.doc
Relative URL: ^/readme.doc
Repository Root: http://svn.red-bean.com/repos/test
Repository UUID: 5e7d134a-54fb-0310-bd04-b611643e5c25
Revision: 1
Node Kind: file
Schedule: normal
Last Changed Author: sally
Last Changed Rev: 42
Last Changed Date: 2003-01-14 23:21:19 -0600 (Tue, 14 Jan 2003)
Lock Token: opaquelocktoken:14011d4b-54fb-0310-8541-dbd16bd471b2
Lock Owner: harry
Lock Created: 2003-01-15 17:35:12 -0600 (Wed, 15 Jan 2003)
Lock Comment (1 line):
My test lock comment
\end{verbatim}

Lastly, \texttt{svn\ info} output is available in XML format by passing
the \texttt{-\/-xml} option:

\begin{verbatim}
$ svn info --xml http://svn.red-bean.com/repos/test
<?xml version="1.0"?>
<info>
<entry
   kind="dir"
   path="."
   revision="1">
<url>http://svn.red-bean.com/repos/test</url>
<relative-url>^/</relative-url>
<repository>
<root>http://svn.red-bean.com/repos/test</root>
<uuid>5e7d134a-54fb-0310-bd04-b611643e5c25</uuid>
</repository>
<wc-info>
<schedule>normal</schedule>
<depth>infinity</depth>
</wc-info>
<commit
   revision="1">
<author>sally</author>
<date>2003-01-15T23:35:12.847647Z</date>
</commit>
</entry>
</info>
\end{verbatim}

\section{svn list (ls)}\label{svn.ref.svn.c.list}

\emph{List directory entries in the repository.}

\textbf{Synopsis}

\texttt{svn\ list\ {[}TARGET{[}@REV{]}...{]}}

\subsection{Description}

List each \textless TARGET\textgreater{} file and the contents of each
\textless TARGET\textgreater{} directory as they exist in the
repository. If \textless TARGET\textgreater{} is a working copy path,
the corresponding repository URL will be used.

The default \textless TARGET\textgreater{} is ``\texttt{.}'', meaning
the repository URL of the current working copy directory.

With \texttt{-\/-verbose} (\texttt{-v}), \texttt{svn\ list} shows the
following fields for each item:

\begin{itemize}
\item
  Revision number of the last commit
\item
  Author of the last commit
\item
  If locked, the letter ``O'' (see the preceding section on
  \hyperref[svn.ref.svn.c.info]{svn info} for details).
\item
  Size (in bytes)
\item
  Date and time of the last commit
\end{itemize}

With \texttt{-\/-xml}, output is in XML format (with a header and an
enclosing document element unless \texttt{-\/-incremental} is also
specified). All of the information is present; the \texttt{-\/-verbose}
(\texttt{-v}) option is not accepted.

\subsection{Options}

\begin{verbatim}
--depth ARG
--incremental
--recursive (-R)
--revision (-r) REV
--verbose (-v)
--xml
\end{verbatim}

\subsection{Examples}

\texttt{svn\ list} is most useful if you want to see what files a
repository has without downloading a working copy:

\begin{verbatim}
$ svn list http://svn.red-bean.com/repos/test/support
README.txt
INSTALL
examples/
…
\end{verbatim}

You can pass the \texttt{-\/-verbose} (\texttt{-v}) option for
additional information, rather like the Unix command \texttt{ls\ -l}:

\begin{verbatim}
$ svn list -v file:///var/svn/repos
     16 sally         28361 Jan 16 23:18 README.txt
     27 sally             0 Jan 18 15:27 INSTALL
     24 harry               Jan 18 11:27 examples/
\end{verbatim}

You can also get \texttt{svn\ list} output in XML format with the
\texttt{-\/-xml} option:

\begin{verbatim}
$ svn list --xml http://svn.red-bean.com/repos/test
<?xml version="1.0"?>
<lists>
<list
   path="http://svn.red-bean.com/repos/test">
<entry
   kind="dir">
<name>examples</name>
<size>0</size>
<commit
   revision="24">
<author>harry</author>
<date>2008-01-18T06:35:53.048870Z</date>
</commit>
</entry>
...
</list>
</lists>
\end{verbatim}

For further details, see the earlier section
\hyperref[svn.tour.history.browsing.list]{Listing versioned
directories}.

\section{svn lock}\label{svn.ref.svn.c.lock}

\emph{Lock working copy paths or URLs in the repository so that no other
user can commit changes to them.}

\textbf{Synopsis}

\texttt{svn\ lock\ TARGET...}

\subsection{Description}

Lock each \textless TARGET\textgreater. If any
\textless TARGET\textgreater{} is already locked by another user, print
a warning and continue locking the rest of the
\textless TARGET\textgreater s. Use \texttt{-\/-force} to steal a lock
from another user or working copy.

\subsection{Options}

\begin{verbatim}
--encoding ENC
--file (-F) FILENAME
--force
--force-log
--message (-m) MESSAGE
--targets FILENAME
\end{verbatim}

\subsection{Examples}

Lock two files in your working copy:

\begin{verbatim}
$ svn lock tree.jpg house.jpg
'tree.jpg' locked by user 'harry'.
'house.jpg' locked by user 'harry'.
\end{verbatim}

Lock a file in your working copy that is currently locked by another
user:

\begin{verbatim}
$ svn lock tree.jpg
svn: warning: W160035: Path '/tree.jpg is already locked by user 'sally' in fi
lesystem '/var/svn/repos/db'
$ svn lock --force tree.jpg
'tree.jpg' locked by user 'harry'.
\end{verbatim}

Lock a file without a working copy:

\begin{verbatim}
$ svn lock http://svn.red-bean.com/repos/test/tree.jpg
'tree.jpg' locked by user 'harry'.
\end{verbatim}

For further details, see \hyperref[svn.advanced.locking]{Locking}.

\section{svn log}\label{svn.ref.svn.c.log}

\emph{Display commit log messages.}

\textbf{Synopsis}

\texttt{svn\ log\ {[}PATH{]}}

\texttt{svn\ log\ URL{[}@REV{]}\ {[}PATH...{]}}

\subsection{Description}

Shows log messages from the repository. If no arguments are supplied,
\texttt{svn\ log} shows the log messages for all files and directories
inside (and including) the current working directory of your working
copy. You can refine the results by specifying a path, one or more
revisions, or any combination of the two. The default revision range for
a local path is \texttt{BASE:1}.

If you specify a URL alone, it prints log messages for everything the
URL contains. If you add paths past the URL, only messages for those
paths under that URL will be printed. The default revision range for a
URL is \texttt{HEAD:1}.

With \texttt{-\/-verbose} (\texttt{-v}), \texttt{svn\ log} will also
print all affected paths with each log message. With \texttt{-\/-quiet}
(\texttt{-q}), \texttt{svn\ log} will not print the log message body
itself, this is compatible with \texttt{-\/-verbose} (\texttt{-v}).

Each log message is printed just once, even if more than one of the
affected paths for that revision were explicitly requested. Logs follow
copy history by default. Use \texttt{-\/-stop-on-copy} to disable this
behavior, which can be useful for determining branch points.

\subsection{Options}

\begin{verbatim}
--change (-c) ARG
--depth ARG
--diff
--diff-cmd CMD
--extensions (-x) ARG
--incremental
--internal-diff
--limit (-l) NUM
--quiet (-q)
--revision (-r) REV
--search ARG
--search-and ARG
--stop-on-copy
--targets FILENAME
--use-merge-history (-g)
--verbose (-v)
--with-all-revprops
--with-no-revprops
--with-revprop ARG
--xml
\end{verbatim}

\subsection{Examples}

You can see the log messages for all the paths that changed in your
working copy by running \texttt{svn\ log} from the top:

\begin{verbatim}
$ svn log
------------------------------------------------------------------------
r20 | harry | 2003-01-17 22:56:19 -0600 (Fri, 17 Jan 2003) | 1 line

Tweak.
------------------------------------------------------------------------
r17 | sally | 2003-01-16 23:21:19 -0600 (Thu, 16 Jan 2003) | 2 lines
…
\end{verbatim}

Examine all log messages for a particular file in your working copy:

\begin{verbatim}
$ svn log foo.c
------------------------------------------------------------------------
r32 | sally | 2003-01-13 00:43:13 -0600 (Mon, 13 Jan 2003) | 1 line

Added defines.
------------------------------------------------------------------------
r28 | sally | 2003-01-07 21:48:33 -0600 (Tue, 07 Jan 2003) | 3 lines
…
\end{verbatim}

If you don\textquotesingle t have a working copy handy, you can log a
URL:

\begin{verbatim}
$ svn log http://svn.red-bean.com/repos/test/foo.c
------------------------------------------------------------------------
r32 | sally | 2003-01-13 00:43:13 -0600 (Mon, 13 Jan 2003) | 1 line

Added defines.
------------------------------------------------------------------------
r28 | sally | 2003-01-07 21:48:33 -0600 (Tue, 07 Jan 2003) | 3 lines
…
\end{verbatim}

If you want several distinct paths underneath the same URL, you can use
the \texttt{URL\ {[}PATH...{]}} syntax:

\begin{verbatim}
$ svn log http://svn.red-bean.com/repos/test/ foo.c bar.c
------------------------------------------------------------------------
r32 | sally | 2003-01-13 00:43:13 -0600 (Mon, 13 Jan 2003) | 1 line

Added defines.
------------------------------------------------------------------------
r31 | harry | 2003-01-10 12:25:08 -0600 (Fri, 10 Jan 2003) | 1 line

Added new file bar.c
------------------------------------------------------------------------
r28 | sally | 2003-01-07 21:48:33 -0600 (Tue, 07 Jan 2003) | 3 lines
…
\end{verbatim}

The \texttt{-\/-verbose} (\texttt{-v}) option causes \texttt{svn\ log}
to include information about the paths that were changed in each
displayed revision. These paths appear, one path per line of output,
with action codes that indicate what type of change was made to the
path.

\begin{verbatim}
$ svn log -v http://svn.red-bean.com/repos/test/ foo.c bar.c
------------------------------------------------------------------------
r32 | sally | 2003-01-13 00:43:13 -0600 (Mon, 13 Jan 2003) | 1 line
Changed paths:
   M /foo.c

Added defines.
------------------------------------------------------------------------
r31 | harry | 2003-01-10 12:25:08 -0600 (Fri, 10 Jan 2003) | 1 line
Changed paths:
   A /bar.c

Added new file bar.c
------------------------------------------------------------------------
r28 | sally | 2003-01-07 21:48:33 -0600 (Tue, 07 Jan 2003) | 3 lines
…
\end{verbatim}

\texttt{svn\ log} uses just a handful of action codes, and they are
similar to the ones the \texttt{svn\ update} command uses:

\begin{description}
\item[\texttt{A}]
The item was added.
\item[\texttt{D}]
The item was deleted.
\item[\texttt{M}]
Properties or textual contents on the item were changed.
\item[\texttt{R}]
The item was replaced by a different one at the same location.
\end{description}

In addition to the action codes which precede the changed paths,
\texttt{svn\ log} with the \texttt{-\/-verbose} (\texttt{-v}) option
will note whether a path was added or replaced as the result of a copy
operation. It does so by printing
\texttt{(from\ COPY-FROM-PATH:COPY-FROM-REV)} after such paths.

When you\textquotesingle re concatenating the results of multiple calls
to the log command, you may want to use the \texttt{-\/-incremental}
option. \texttt{svn\ log} normally prints out a dashed line at the
beginning of a log message, after each subsequent log message, and
following the final log message. If you ran \texttt{svn\ log} on a range
of two revisions, you would get this:

\begin{verbatim}
$ svn log -r 14:15
------------------------------------------------------------------------
r14 | …

------------------------------------------------------------------------
r15 | …

------------------------------------------------------------------------
\end{verbatim}

However, if you wanted to gather two nonsequential log messages into a
file, you might do something like this:

\begin{verbatim}
$ svn log -r 14 > mylog
$ svn log -r 19 >> mylog
$ svn log -r 27 >> mylog
$ cat mylog
------------------------------------------------------------------------
r14 | …

------------------------------------------------------------------------
------------------------------------------------------------------------
r19 | …

------------------------------------------------------------------------
------------------------------------------------------------------------
r27 | …

------------------------------------------------------------------------
\end{verbatim}

You can avoid the clutter of the double dashed lines in your output by
using the \texttt{-\/-incremental} option:

\begin{verbatim}
$ svn log --incremental -r 14 > mylog
$ svn log --incremental -r 19 >> mylog
$ svn log --incremental -r 27 >> mylog
$ cat mylog
------------------------------------------------------------------------
r14 | …

------------------------------------------------------------------------
r19 | …

------------------------------------------------------------------------
r27 | …
\end{verbatim}

The \texttt{-\/-incremental} option provides similar output control when
using the \texttt{-\/-xml} option:

\begin{verbatim}
$ svn log --xml --incremental -r 1 sandwich.txt
<logentry
   revision="1">
<author>harry</author>
<date>2008-06-03T06:35:53.048870Z</date>
<msg>Initial Import.</msg>
</logentry>
\end{verbatim}

Sometimes when you run \texttt{svn\ log} on a specific path and a
specific revision, you see no log information output at all, as in the
following:

\begin{verbatim}
$ svn log -r 20 http://svn.red-bean.com/untouched.txt
------------------------------------------------------------------------
\end{verbatim}

That just means the path wasn\textquotesingle t modified in that
revision. To get log information for that revision, either run the log
operation against the repository\textquotesingle s root URL, or specify
a path that you happen to know was changed in that revision:

\begin{verbatim}
$ svn log -r 20 touched.txt 
------------------------------------------------------------------------
r20 | sally | 2003-01-17 22:56:19 -0600 (Fri, 17 Jan 2003) | 1 line

Made a change.
------------------------------------------------------------------------
\end{verbatim}

Beginning with Subversion 1.7, users can take advantage of a special
output mode which combines the information from \texttt{svn\ log} with
what you would see when running \texttt{svn\ diff} on the same targets
for each revision of the log. Simply invoke \texttt{svn\ log} with the
\texttt{-\/-diff} option to trigger this output mode.

\begin{verbatim}
$ svn log -r 20 touched.txt --diff
------------------------------------------------------------------------
r20 | sally | 2003-01-17 22:56:19 -0600 (Fri, 17 Jan 2003) | 1 line

Made a change.

Index: touched.txt
===================================================================
--- touched.txt (revision 19)
+++ touched.txt (revision 20)
@@ -1 +1,2 @@
 This is the file 'touched.txt'.
+We add such exciting text to files around here!
------------------------------------------------------------------------
$
\end{verbatim}

As with \texttt{svn\ diff}, you may also make use of many of the various
options which control the way the difference is generated, including
\texttt{-\/-depth}, \texttt{-\/-diff-cmd}, and \texttt{-\/-extensions}
(\texttt{-x}).

Beginning with Subversion 1.8, users can filter \texttt{svn\ log} output
using \texttt{-\/-search} and \texttt{-\/-search-and} options. When
using these options, a log message is shown only if a
revision\textquotesingle s author, date, log message text, or list of
changed paths, matches a search pattern. Searching by changed patch
requires \texttt{-\/-verbose} option, otherwise \texttt{svn\ log} does
not show changed paths therefore they can\textquotesingle t be filtered.

The search pattern (also called glob or shell wildcard pattern) can
contain regular characters and the following wildcard characters:

\begin{description}
\item[\texttt{?}]
Matches any single character.
\item[\texttt{*}]
Matches a sequence of arbitrary characters.
\item[\texttt{{[}ABC{]}}]
Matches any of the characters listed inside the brackets.
\end{description}

Using multiple \texttt{-\/-search} parameters will show log messages
that match the pattern specified at least in one of the options. For
example:

\begin{verbatim}
$ svn log --search sally --search harry https://svn.red-bean.com/repos/test
------------------------------------------------------------------------
r1701 | sally | 2011-10-12 22:35:30 -0600 (Wed, 12 Oct 2011) | 1 line

Add a reminder.
------------------------------------------------------------------------
r1564 | harry | 2011-10-09 22:35:30 -0600 (Sun, 09 Oct 2011) | 1 line

Merge r1560 to the 1.0.x branch.
------------------------------------------------------------------------
$
        
\end{verbatim}

Using \texttt{-\/-search} with \texttt{-\/-search-and} options will show
log messages that match the combined pattern from both options. For
example:

\begin{verbatim}
$ svn log --verbose --search sally --search-and /foo/bar https://svn.red-bean.com/repos/test
------------------------------------------------------------------------
r1555 | sally | 2011-07-15 22:33:14 -0600 (Fri, 15 Jul 2011) | 1 line
Changed paths:
M /foo/bar/src.c

Typofix.
------------------------------------------------------------------------
r1530 | sally | 2011-07-13 07:24:11 -0600 (Wed, 13 Jul 2011) | 1 line
Changed paths:
M /foo/bar
M /foo/build

Fix up some svn:ignore properties.
------------------------------------------------------------------------
$
\end{verbatim}

\texttt{-\/-search} and \texttt{-\/-search-and} options does not
actually perform a search. They just filter the \texttt{svn\ log} output
to display only log messages that match the specified pattern.
Therefore, if \texttt{-\/-limit} is used, it restricts the number of log
messages searched, rather than restricting the output to a particular
number of matching log messages.

\section{svn merge}\label{svn.ref.svn.c.merge}

\emph{Apply the differences between two sources to a working copy path.}

\textbf{Synopsis}

\texttt{svn\ merge\ SOURCE{[}@REV{]}\ {[}TARGET\_WCPATH{]}}

\texttt{svn\ merge\ {[}-c\ M{[},N...{]}\ \textbar{}\ -r\ N:M\ ...{]}\ SOURCE{[}@REV{]}\ {[}TARGET\_WCPATH{]}}

\texttt{svn\ merge\ SOURCE1{[}@N{]}\ SOURCE2{[}@M{]}\ {[}TARGET\_WCPATH{]}}

\subsection{Description}

In all three forms \textless TARGET\_WCPATH\textgreater{} is the working
copy path that will receive the differences. If
\textless TARGET\_WCPATH\textgreater{} is omitted, the changes are
applied to the current working directory, unless the sources have
identical basenames that match a file within the current working
directory. In that case, the differences will be applied to that file.

In the first two forms, \textless SOURCE\textgreater{} can be either a
URL or a working copy path (in which case its corresponding URL is
used). If the peg revision \textless REV\textgreater{} is not specified,
then \texttt{HEAD} is assumed. In the third form the same rules apply
for \textless SOURCE1\textgreater, \textless SOURCE2\textgreater,
\textless M\textgreater, and \textless N\textgreater{} with the only
difference being that if either source is a working copy path, then the
peg revisions \emph{must} be explicitly stated.

\begin{itemize}
\item
  Automatic Merges

  The first form is called an ``automatic merge'' and is used to perform
  ``sync'' and ``reintegrate'' merges. ``Sync'' merges merge eligible
  changes to a branch (\textless TARGET\_WCPATH\textgreater) from the
  branch\textquotesingle s ancestor branch
  (\textless SOURCE\textgreater). ``Eligible'' changes are defined as
  those that were not previously merged from
  \textless SOURCE\textgreater{} to
  \textless TARGET\_WCPATH\textgreater. See
  \hyperref[svn.branchmerge.basicmerging.stayinsync]{Keeping a Branch in
  Sync}. ``Reintegrate'' merges merge changes from a feature branch
  (\textless SOURCE\textgreater) back into the feature
  branch\textquotesingle s ancestor branch
  (\textless TARGET\_WCPATH\textgreater), see
  \hyperref[svn.branchmerge.basicmerging.reintegrate]{Reintegrating a
  Branch} and \hyperref[svn.branchmerge.commonpatterns.feature]{Feature
  Branches}.
\item
  Cherrypick Merges

  The second form is called a ``cherry-pick'' merge and is used to merge
  an explicitly defined set of changes from one branch to another.
  \textless SOURCE\textgreater{} in revision \textless REV\textgreater{}
  is compared as it existed between revisions \textless N\textgreater{}
  and \textless M\textgreater{} for each revision range provided. See
  \hyperref[svn.branchmerge.cherrypicking]{Cherrypicking} for more
  information.

  Multiple \texttt{-c} and/or \texttt{-r} instances may be specified,
  and mixing of forward and reverse ranges is allowed--- the ranges are
  internally compacted to their minimum representation before merging
  begins (which may result in a no-op merge or conflicts that cause the
  merge to stop before merging all of the requested revisions).
\item
  2-URL Merges

  In the third form, called a ``2-URL Merge'', the difference between
  \textless SOURCE1\textgreater{} at revision \textless N\textgreater{}
  and \textless SOURCE2\textgreater{} at revision
  \textless M\textgreater{} is generated and applied to
  \textless TARGET\_WCPATH\textgreater. The revisions default to
  \texttt{HEAD} if omitted.
\end{itemize}

If \hyperref[svn.branchmerge.basicmerging.mergetracking]{Merge Tracking}
is active, then Subversion will internally track metadata (i.e. the
\texttt{svn:mergeinfo} property) about merge operations when the two
merge sources are ancestrally related---if the first source is an
ancestor of the second or vice versa---this is guaranteed to be the case
when performing automatic merges. Subversion will also take preexisting
merge metadata on the working copy target into account when determining
what revisions to merge and in an effort to avoid repeat merges and
needless conflicts it may only merge a subset of the requested ranges.

Unlike \texttt{svn\ diff}, the merge command takes the ancestry of a
file into consideration when performing a merge operation. This is very
important when you\textquotesingle re merging changes from one branch
into another and you\textquotesingle ve renamed a file on one branch but
not the other.

The \texttt{-\/-ignore-ancestry} option will cause
\hyperref[svn.branchmerge.basicmerging.mergetracking]{Merge Tracking} to
be disabled and makes merge act like \texttt{svn\ diff}, ignoring the
ancestry of files when merging.

\subsection{Options}

\begin{verbatim}
--accept ACTION
--allow-mixed-revisions
--change (-c) ARG
--depth ARG
--diff3-cmd CMD
--dry-run
--extensions (-x) ARG
--force
--ignore-ancestry
--quiet (-q)
--record-only
--reintegrate
--revision (-r) REV
--verbose (-v)
\end{verbatim}

\subsection{Examples}

Reintegrate a branch back into the trunk---assuming that you have an
up-to-date working copy of the trunk (the \texttt{-\/-verbose} option
prints additional information regarding what the merge is doing prior to
actually applying any diff; useful in very large which might take a
significant amount of time to complete):

\begin{verbatim}
$ svn merge ^/branches/feature-branch-calc-enhancements trunk --verbose
checking branch relationship...
calculating automatic merge...
merging...
--- Merging r12 through r37 into 'trunk':
U    trunk/calc/brush.c
--- Recording mergeinfo for merge of r12 through r37 into 'trunk':
 U   trunk

$ # build, test, verify, ...

$ svn commit trunk -m "Reintegrate the calc enhancements back to trunk!"
Sending        trunk
Sending        trunk/calc/brush.c
Transmitting file data .
Committed revision 38.
\end{verbatim}

Cherry-pick merge a single change to a file:

\begin{verbatim}
$ svn merge ^/trunk/calc/brush.c branches/1.x/calc/brush.c -c38
--- Merging r38 into 'branches/1.x/calc/brush.c':
U    branches/1.x/calc/brush.c
--- Recording mergeinfo for merge of r38 into 'branches/1.x/calc/brush.c':
 G   branches/1.x/calc/brush.c
\end{verbatim}

Merge the differences between two unrelated branches into a third
branch:

\begin{verbatim}
$ svn merge ^/vendor-drop/vendor-1.0 ^/vendor-drop/vendor-1.1 \
            trunk --ignore-ancestry
--- Merging differences between repository URLs into 'trunk':
U    trunk/draw/draw.py
\end{verbatim}

\section{svn mergeinfo}\label{svn.ref.svn.c.mergeinfo}

\emph{Query merge-related information. See
\hyperref[svn.branchmerge.basicmerging.mergeinfo]{Mergeinfo and
Previews} for details.}

\textbf{Synopsis}

\texttt{svn\ mergeinfo\ SOURCE\_URL{[}@REV{]}\ {[}TARGET{[}@REV{]}{]}}

\subsection{Description}

Query information related to merges (or potential merges) between
\textless SOURCE-URL\textgreater{} and \textless TARGET\textgreater. If
the \texttt{-\/-show-revs} option is not provided, display a graphical
representation of revisions which have been fully merged from
\textless SOURCE-URL\textgreater{} to \textless TARGET\textgreater.
Otherwise, list either the \texttt{merged} or \texttt{eligible}
revisions as specified by the \texttt{-\/-show-revs} option.

\subsection{Options}

\begin{verbatim}
--depth ARG
--recursive (-R)
--revision (-r) REV
--show-revs ARG
\end{verbatim}

\subsection{Examples}

Graphical summary of the merges from one branch to another:

\begin{verbatim}
$ svn mergeinfo ^/trunk feature-branch
    youngest  last               repos.
    common    full     tip of    path of
    ancestor  merge    branch    branch

    11        16       33
    |         |        |
  -------| |------------         trunk
     \         \
      \         \
       --| |------------         feature-branch
                       |
                       33
\end{verbatim}

List the operative revisions merged from one branch to another:

\begin{verbatim}
$ svn mergeinfo ^/trunk feature-branch --show-revs merged
r15
r16
\end{verbatim}

List the operative revisions eligible to be merged from one branch to
another:

\begin{verbatim}
$ svn mergeinfo ^/trunk feature-branch --show-revs eligible
r28
r30
\end{verbatim}

\section{svn mkdir}\label{svn.ref.svn.c.mkdir}

\emph{Create a new directory under version control.}

\textbf{Synopsis}

\texttt{svn\ mkdir\ PATH...}

\texttt{svn\ mkdir\ URL...}

\subsection{Description}

Create a directory with a name given by the final component of the
\textless PATH\textgreater{} or \textless URL\textgreater. A directory
specified by a working copy \textless PATH\textgreater{} is scheduled
for addition in the working copy. A directory specified by a URL is
created in the repository via an immediate commit. Multiple directory
URLs are committed atomically. In both cases, all the intermediate
directories must already exist unless the \texttt{-\/-parents} option is
used.

\subsection{Options}

\begin{verbatim}
--editor-cmd CMD
--encoding ENC
--file (-F) FILENAME
--force-log
--message (-m) MESSAGE
--parents
--quiet (-q)
--with-revprop ARG
\end{verbatim}

\subsection{Examples}

Create a directory in your working copy:

\begin{verbatim}
$ svn mkdir newdir
A         newdir
\end{verbatim}

Create one in the repository (this is an instant commit, so a log
message is required):

\begin{verbatim}
$ svn mkdir -m "Making a new dir." http://svn.red-bean.com/repos/newdir

Committed revision 26.
\end{verbatim}

\section{svn move (mv)}\label{svn.ref.svn.c.move}

\emph{Move a file or directory.}

\textbf{Synopsis}

\texttt{svn\ move\ SRC...\ DST}

\subsection{Description}

This command moves files or directories in your working copy or in the
repository.

This command is equivalent to an \texttt{svn\ copy} followed by
\texttt{svn\ delete}.

When moving multiple sources, they will be added as children of
\textless DST\textgreater, which must be a directory.

Subversion does not support moving between working copies and URLs. In
addition, you can only move files within a single
repository---Subversion does not support cross-repository moving.
Subversion supports the following types of moves within a single
repository:

\begin{description}
\item[WC → WC]
Move and schedule a file or directory for addition (with history).
\item[URL → URL]
Complete server-side rename.
\end{description}

When moving large trees you should be aware that the URL → URL moves are
lighter than WC → WC moves. Moving nodes inside a working copy does more
than just change directory listings (it will copy files, manage
temporary files, and expand keywords) and may be significantly slower.

Also bear in mind that a WC → WC move in a mixed-revision working copy
may yield unexpected results (see
\hyperref[svn.basic.in-action.mixedrevs]{Mixed-revision working
copies}).

\subsection{Options}

\begin{verbatim}
--editor-cmd CMD
--encoding ENC
--file (-F) FILENAME
--force
--force-log
--message (-m) MESSAGE
--parents
--quiet (-q)
--revision (-r) REV
--with-revprop ARG
\end{verbatim}

\subsection{Examples}

Move a file in your working copy:

\begin{verbatim}
$ svn move foo.c bar.c
A         bar.c
D         foo.c
\end{verbatim}

Move several files in your working copy into a subdirectory:

\begin{verbatim}
$ svn move baz.c bat.c qux.c src
A         src/baz.c
D         baz.c
A         src/bat.c
D         bat.c
A         src/qux.c
D         qux.c
\end{verbatim}

Move a file in the repository (this is an immediate commit, so it
requires a commit message):

\begin{verbatim}
$ svn move -m "Move a file" http://svn.red-bean.com/repos/foo.c \
                            http://svn.red-bean.com/repos/bar.c

Committed revision 27.
\end{verbatim}

\section{svn patch}\label{svn.ref.svn.c.patch}

\emph{Apply changes represented in a unidiff patch to the working copy.}

\textbf{Synopsis}

\texttt{svn\ patch\ PATCHFILE\ {[}WCPATH{]}}

\subsection{Description}

This subcommand will apply changes described a unidiff-formatted patch
file \textless PATCHFILE\textgreater{} to the working copy
\textless WCPATH\textgreater. As with most other working copy
subcommands, if \textless WCPATH\textgreater{} is omitted, the changes
are applied to the current working directory. A unidiff patch suitable
for application to a working copy can be produced with the
\texttt{svn\ diff} command or third-party differencing tools. Any
non-unidiff content found in the patch file is ignored.

Changes listed in the patch file will either be applied or rejected. If
a change does not match at its exact line offset, it may be applied
earlier or later in the file if a match is found elsewhere for the
surrounding lines of context provided by the patch. A change may also be
applied with fuzz---meaning, one or more lines of context are ignored
when attempting to match the change location. If no matching context can
be found for a change, the change conflicts and will be written to a
reject file which bears the extension \texttt{.svnpatch.rej}.

\texttt{svn\ patch} reports a status line for patched file or directory
using letter codes, very similar to the way that \texttt{svn\ update}
provides notification. The letter codes have the following meanings:

\begin{description}
\item[\texttt{A}]
Added
\item[\texttt{D}]
Deleted
\item[\texttt{C}]
Conflicted
\item[\texttt{G}]
Merged
\item[\texttt{U}]
Updated
\end{description}

Changes applied with an offset or fuzz are reported on lines starting
with the \textquotesingle{}\texttt{\textgreater{}}\textquotesingle{}
symbol. You should review such changes carefully.

If the patch removes all content from a file, that file is automatically
scheduled for deletion. Likewise, if the patch creates a new file, that
file is automatically scheduled for addition. Use \texttt{svn\ revert}
to undo undesired deletions and additions.

\subsection{Options}

\begin{verbatim}
--dry-run
--ignore-whitespace
--quiet (-q)
--reverse-diff
--strip NUM
\end{verbatim}

\subsection{Examples}

Apply a simple patch file generated by the \texttt{svn\ diff} command.
Our patch file will create a new file, delete another file, and modify a
third\textquotesingle s contents and properties. Here\textquotesingle s
the patch file itself (which we\textquotesingle ll assume is creatively
named \texttt{PATCH}):

\begin{verbatim}
Index: deleted-file
===================================================================
--- deleted-file    (revision 3)
+++ deleted-file    (working copy)
@@ -1 +0,0 @@
-This file will be deleted.
Index: changed-file
===================================================================
--- changed-file    (revision 4)
+++ changed-file    (working copy)
@@ -1,6 +1,6 @@
 The letters in a line of text
 Could make your day much better.
 But expanded into paragraphs,
-I'd tell of kangaroos and calves
+I'd tell of monkeys and giraffes
 Until you were all smiles and laughs
 From my letter made of letters.

Property changes on: changed-file
___________________________________________________________________
Added: propname
## -0,0 +1 ##
+propvalue
Index: added-file
===================================================================
--- added-file  (revision 0)
+++ added-file  (working copy)
@@ -0,0 +1 @@
+This is an added file.
\end{verbatim}

We can apply the previous patch file to another working copy from our
repository using \texttt{svn\ patch}, and verify that it did the right
thing by using \texttt{svn\ diff}:

\begin{verbatim}
$ cd /some/other/workingcopy
$ svn patch /path/to/PATCH
D         deleted-file
UU        changed-file
A         added-file
$ svn diff
Index: deleted-file
===================================================================
--- deleted-file    (revision 3)
+++ deleted-file    (working copy)
@@ -1 +0,0 @@
-This file will be deleted.
Index: changed-file
===================================================================
--- changed-file    (revision 4)
+++ changed-file    (working copy)
@@ -1,6 +1,6 @@
 The letters in a line of text
 Could make your day much better.
 But expanded into paragraphs,
-I'd tell of kangaroos and calves
+I'd tell of monkeys and giraffes
 Until you were all smiles and laughs
 From my letter made of letters.

Property changes on: changed-file
___________________________________________________________________
Added: propname
## -0,0 +1 ##
+propvalue
Index: added-file
===================================================================
--- added-file  (revision 0)
+++ added-file  (working copy)
@@ -0,0 +1 @@
+This is an added file.
$
\end{verbatim}

Sometimes you might need Subversion to interpret a patch ``in
reverse''---where added things get treated as removed things, and
vice-versa. Use the \texttt{-\/-reverse-diff} option for this purpose.
In the following example, we\textquotesingle ll squirrel away a patch
file which describes the changes in our working copy, and then use a
reverse patch operation to undo those changes.

\begin{verbatim}
$ svn status
M       foo.c
$ svn diff > PATCH
$ cat PATCH
Index: foo.c
===================================================================
--- foo.c   (revision 128)
+++ foo.c   (working copy)
@@ -1003,7 +1003,7 @@
     return ERROR_ON_THE_G_STRING;
 
   /* Do something in a loop. */
-  for (i = 0; i < txns->nelts; i++)
+  for (i = 0; i < txns->nelts; i--)
     {
       status = do_something(i);
       if (status)
$ svn patch --reverse-diff PATCH
U         foo.c
$ svn status
$
\end{verbatim}

\section{svn propdel (pdel, pd)}\label{svn.ref.svn.c.propdel}

\emph{Remove a property from an item.}

\textbf{Synopsis}

\texttt{svn\ propdel\ PROPNAME\ {[}PATH...{]}}

\texttt{svn\ propdel\ PROPNAME\ -\/-revprop\ -r\ REV\ {[}TARGET{]}}

\subsection{Description}

This removes properties from files, directories, or revisions. The first
form removes versioned properties in your working copy, and the second
removes unversioned remote properties on a repository revision
(\textless TARGET\textgreater{} determines only which repository to
access).

\subsection{Options}

\begin{verbatim}
--changelist (--cl) ARG
--depth ARG
--quiet (-q)
--recursive (-R)
--revision (-r) REV
--revprop
\end{verbatim}

\subsection{Examples}

Delete a property from a file in your working copy:

\begin{verbatim}
$ svn propdel svn:mime-type some-script
property 'svn:mime-type' deleted from 'some-script'.
\end{verbatim}

Delete a revision property:

\begin{verbatim}
$ svn propdel --revprop -r 26 release-date 
property 'release-date' deleted from repository revision '26'
\end{verbatim}

\section{svn propedit (pedit, pe)}\label{svn.ref.svn.c.propedit}

\emph{Edit the property of one or more items under version control. See
\hyperref[svn.ref.svn.c.propset]{svn propset (pset, ps)} later in this
chapter.}

\textbf{Synopsis}

\texttt{svn\ propedit\ PROPNAME\ TARGET...}

\texttt{svn\ propedit\ PROPNAME\ -\/-revprop\ -r\ REV\ {[}TARGET{]}}

\subsection{Description}

Edit one or more properties using your favorite editor. The first form
edits versioned properties in your working copy, and the second edits
unversioned remote properties on a repository revision
(\textless TARGET\textgreater{} determines only which repository to
access).

\subsection{Options}

\begin{verbatim}
--editor-cmd CMD
--encoding ENC
--file (-F) FILENAME
--force
--force-log
--message (-m) MESSAGE
--revision (-r) REV
--revprop
--with-revprop ARG
\end{verbatim}

\subsection{Examples}

\texttt{svn\ propedit} makes it easy to modify properties that have
multiple values:

\begin{verbatim}
$ svn propedit svn:keywords foo.c 

    # svn will open in your favorite text editor a temporary file
    # containing the current contents of the svn:keywords property.  You
    # can add multiple values to a property easily here by entering one
    # value per line.  When you save the temporary file and exit,
    # Subversion will re-read the temporary file and use its updated
    # contents as the new value of the property.

Set new value for property 'svn:keywords' on 'foo.c'
$
\end{verbatim}

\section{svn propget (pget, pg)}\label{svn.ref.svn.c.propget}

\emph{Print the value of a property.}

\textbf{Synopsis}

\texttt{svn\ propget\ PROPNAME\ {[}TARGET{[}@REV{]}...{]}}

\texttt{svn\ propget\ PROPNAME\ -\/-revprop\ -r\ REV\ {[}URL{]}}

\subsection{Description}

Print the value of a property on files, directories, or revisions. The
first form prints the versioned property of an item or items in your
working copy, and the second prints unversioned remote properties on a
repository revision. See \hyperref[svn.advanced.props]{Properties} for
more information on properties.

\subsection{Options}

\begin{verbatim}
--changelist (--cl) ARG
--depth ARG
--recursive (-R)
--revision (-r) REV
--revprop
--show-inherited-props
--strict
--verbose (-v)
--xml
\end{verbatim}

\subsection{Examples}

Examine a property of a file in your working copy:

\begin{verbatim}
$ svn propget svn:keywords foo.c
Author
Date
Rev
\end{verbatim}

The same goes for a revision property:

\begin{verbatim}
$ svn propget svn:log --revprop -r 20 
Began journal.
\end{verbatim}

For a more structured display of properties, use the
\texttt{-\/-verbose} (\texttt{-v}) option:

\begin{verbatim}
$ svn propget svn:keywords foo.c --verbose
Properties on 'foo.c':
  svn:keywords
    Author
    Date
    Rev
\end{verbatim}

Examine the versioned properties inherited by a URL in your repository
using the \texttt{-\/-show-inherited-props} option:

\begin{verbatim}
$ svn pg svn:global-ignores --verbose --show-inherited-props ^/branches/1.x
Inherited properties on 'http://svn.example.com/repos/branches/1.x',
from 'http://svn.example.com/repos':
  svn:global-ignores
    *.diff
    *.patch
\end{verbatim}

By default, \texttt{svn\ propget} will append a trailing end-of-line
sequence to the property value it prints. Most of the time, this is a
desirable feature that has a positive effect on the printed output. But
there are times when you might wish to capture the precise property
value, perhaps because that value is not textual in nature, but of some
binary format (such as a JPEG thumbnail stored as a property value, for
example). To disable pretty-printing of property values, use the
\texttt{-\/-strict} option.

Lastly, you can get \texttt{svn\ propget} output in XML format with the
\texttt{-\/-xml} option:

\begin{verbatim}
$ svn propget --xml svn:ignore .
<?xml version="1.0"?>
<properties>
<target
   path="">
<property
   name="svn:ignore">*.o
</property>
</target>
</properties>
\end{verbatim}

\section{svn proplist (plist, pl)}\label{svn.ref.svn.c.proplist}

\emph{List all properties.}

\textbf{Synopsis}

\texttt{svn\ proplist\ {[}TARGET{[}@REV{]}...{]}}

\texttt{svn\ proplist\ -\/-revprop\ -r\ REV\ {[}TARGET{]}}

\subsection{Description}

List all properties on files, directories, or revisions. The first form
lists versioned properties in your working copy, and the second lists
unversioned remote properties on a repository revision
(\textless TARGET\textgreater{} determines only which repository to
access).

\subsection{Options}

\begin{verbatim}
--changelist (--cl) ARG
--depth ARG
--quiet (-q)
--recursive (-R)
--revision (-r) REV
--revprop
--show-inherited-props
--verbose (-v)
--xml
\end{verbatim}

\subsection{Examples}

You can use \texttt{proplist} to see the properties on an item in your
working copy:

\begin{verbatim}
$ svn proplist foo.c
Properties on 'foo.c':
  svn:mime-type
  svn:keywords
  owner
\end{verbatim}

But with the \texttt{-\/-verbose} (\texttt{-v}) flag,
\texttt{svn\ proplist} is extremely handy as it also shows you the
values for the properties:

\begin{verbatim}
$ svn proplist -v foo.c
Properties on 'foo.c':
  svn:mime-type
    text/plain
  svn:keywords
    Author Date Rev
  owner
    sally
\end{verbatim}

List all the versioned properties inherited by a file in your working
copy using the \texttt{-\/-show-inherited-props} option:

\begin{verbatim}
$ svn proplist --show-inherited-props foo.c
Inherited properties on 'foo.c',
from 'http://svn.example.com/repos':
  svn:auto-props
  svn:global-ignores
Inherited properties on 'foo.c',
from '/home/theob/svn/working-copies/baz-wc':
  svn:auto-props
\end{verbatim}

Lastly, you can get \texttt{svn\ proplist} output in XML format with the
\texttt{-\/-xml} option:

\begin{verbatim}
$ svn proplist --xml 
<?xml version="1.0"?>
<properties>
<target
   path=".">
<property
   name="svn:ignore"/>
</target>
</properties>
\end{verbatim}

\section{svn propset (pset, ps)}\label{svn.ref.svn.c.propset}

\emph{Set \textless PROPNAME\textgreater{} to
\textless PROPVAL\textgreater{} on files, directories, or revisions.}

\textbf{Synopsis}

\texttt{svn\ propset\ PROPNAME\ {[}PROPVAL\ \textbar{}\ -F\ VALFILE{]}\ PATH...}

\texttt{svn\ propset\ PROPNAME\ -\/-revprop\ -r\ REV\ {[}PROPVAL\ \textbar{}\ -F\ VALFILE{]}\ {[}TARGET{]}}

\subsection{Description}

Set \textless PROPNAME\textgreater{} to \textless PROPVAL\textgreater{}
on files, directories, or revisions. The first example creates a
versioned, local property change in the working copy, and the second
creates an unversioned, remote property change on a repository revision
(\textless TARGET\textgreater{} determines only which repository to
access).

Subversion has a number of ``special'' properties that affect its
behavior. See
\hyperref[svn.advanced.props.ref]{Subversion\textquotesingle s Reserved
Properties} for details.

\subsection{Options}

\begin{verbatim}
--changelist (--cl) ARG
--depth ARG
--encoding ENC
--file (-F) FILENAME
--force
--quiet (-q)
--recursive (-R)
--revision (-r) REV
--revprop
--targets FILENAME
\end{verbatim}

\subsection{Examples}

Set the MIME type for a file:

\begin{verbatim}
$ svn propset svn:mime-type image/jpeg foo.jpg 
property 'svn:mime-type' set on 'foo.jpg'
\end{verbatim}

On a Unix system, if you want a file to have the executable permission
set:

\begin{verbatim}
$ svn propset svn:executable ON somescript
property 'svn:executable' set on 'somescript'
\end{verbatim}

Perhaps you have an internal policy to set certain properties for the
benefit of your coworkers:

\begin{verbatim}
$ svn propset owner sally foo.c
property 'owner' set on 'foo.c'
\end{verbatim}

If you made a mistake in a log message for a particular revision and
want to change it, use \texttt{-\/-revprop} and set \texttt{svn:log} to
the new log message:

\begin{verbatim}
$ svn propset --revprop -r 25 svn:log "Journaled about trip to New York."
property 'svn:log' set on repository revision '25'
\end{verbatim}

Or, if you don\textquotesingle t have a working copy, you can provide a
URL:

\begin{verbatim}
$ svn propset --revprop -r 26 svn:log "Document nap." \
              http://svn.red-bean.com/repos
property 'svn:log' set on repository revision '25'
\end{verbatim}

Lastly, you can tell \texttt{propset} to take its input from a file. You
could even use this to set the contents of a property to something
binary:

\begin{verbatim}
$ svn propset owner-pic -F sally.jpg moo.c 
property 'owner-pic' set on 'moo.c'
\end{verbatim}

By default, you cannot modify revision properties in a Subversion
repository. Your repository administrator must explicitly enable
revision property modifications by creating a hook named
\texttt{pre-revprop-change}. See
\hyperref[svn.reposadmin.hooks]{Implementing Repository Hooks} for more
information on hook scripts.

\section{svn relocate}\label{svn.ref.svn.c.relocate}

\emph{Relocate the working copy to point to a different repository root
URL.}

\textbf{Synopsis}

\texttt{svn\ relocate\ FROM-PREFIX\ TO-PREFIX\ {[}PATH...{]}}

\texttt{svn\ relocate\ TO-URL\ {[}PATH{]}}

\subsection{Description}

Sometimes an administrator might change the location (or apparent
location, from the client\textquotesingle s perspective) of a
repository. The content of the repository doesn\textquotesingle t
change, but the repository\textquotesingle s root URL does. The hostname
may change because the repository is now being served from a different
computer. Or, perhaps the URL scheme changes because the repository is
now being served via SSL (using \texttt{https://}) instead of over plain
HTTP. There are many different reasons for these types of repository
relocations. But ideally, a ``change of address'' for a repository
shouldn\textquotesingle t suddently cause all the working copies which
point to that repository to become forever unusable. And fortunately,
that\textquotesingle s not the case. Rather than force users to check
out a new working copy when a repository is relocated, Subversion
provides the \texttt{svn\ relocate} command, which ``rewrites'' the
working copy\textquotesingle s administrative metadata to refer to the
new repository location.

The first \texttt{svn\ relocate} syntax allows you to update one or more
working copies by what essentially amounts to a find-and-replace within
the repository root URLs recorded in those working copies. Subversion
will replace the initial substring \textless FROM-PREFIX\textgreater{}
with the string \textless TO-PREFIX\textgreater{} in those URLs. These
initial URL substrings can be as long or as short as is necessary to
differentiate between them. Obviously, to use this syntax form, you need
to know both the current root URL of the repository to which the working
copy is pointing, and the new URL of that repository. (You can use
\texttt{svn\ info} to determine the former.)

The second syntax does not require that you know the current repository
root URL with which the working copy is associated at all---only the new
repository URL (\textless TO-URL\textgreater) to which it should be
pointing. In this syntax form, only one working copy may be relocated at
a time.

Users often get confused about the difference between
\texttt{svn\ switch} and \texttt{svn\ relocate}. Here\textquotesingle s
the rule of thumb:

\begin{itemize}
\item
  If the working copy needs to reflect a new directory \emph{within} the
  repository, use \texttt{svn\ switch}.
\item
  If the working copy still reflects the same repository directory, but
  the location of the repository itself has changed, use
  \texttt{svn\ relocate}.
\end{itemize}

\subsection{Options}

\begin{verbatim}
--ignore-externals
\end{verbatim}

\subsection{Examples}

Let\textquotesingle s start with a working copy that reflects a local
repository URL:

\begin{verbatim}
$ svn info | grep ^URL:
URL: file:///var/svn/repos/trunk
$
\end{verbatim}

One day the administrator decides to rename the on-disk repository
directory. We missed the memo, so we see an error the next time we try
to update our working copy.

\begin{verbatim}
$ svn up
Updating '.':
svn: E180001: Unable to connect to a repository at URL 'file:///var/svn/repos/trunk'
\end{verbatim}

After cornering the administrator over by the vending machines, we learn
about the repository being moved and are told the new URL. Rather than
checkout a new working copy, though, we simply ask Subversion to rewrite
the working copy metadata to point to the new repository location.

\begin{verbatim}
$ svn relocate file:///var/svn/new-repos/trunk
$
\end{verbatim}

Subversion doesn\textquotesingle t tell us much about what it did, but
hey---error-free operation is really all we need, right? Our working
copy is functional for online operations again.

\begin{verbatim}
$ svn up
Updating '.':
A    lib/new.c
M    src/code.h
M    src/headers.h
…
\end{verbatim}

Once again, this type of relocation affects \emph{working copy metadata
only}. It will not change your versioned or unversioned file contents,
perform any version control operations (such as commits or updates), and
so on.

A few months later, we\textquotesingle re told that the company is
moving development to separate machines and that we\textquotesingle ll
be using HTTP to access the repository. So we relocate our working copy
again.

\begin{verbatim}
$ svn relocate http://svn.company.com/repos/trunk
$
\end{verbatim}

Now, each time we perform a relocation of this sort, Subversion contacts
the repository---at its new URL, of course---to verify a few things.

First, it wants to compare the UUID of the repository against what is
stored in the working copy. If these UUIDs don\textquotesingle t match,
the working copy relocation is disallowed. Maybe this
isn\textquotesingle t the same repository (just in a new location) after
all?

Secondly, Subversion wants to ensure that the updated working copy
metadata jives with respect to the directory location \emph{inside} the
repository. Subversion won\textquotesingle t let you accidentally
relocate a working copy of a branch in your repository to the URL of a
different branch in the same repository. (That\textquotesingle s what
\texttt{svn\ switch}, described in \hyperref[svn.ref.svn.c.switch]{svn
switch (sw)}, is for.)

Also, Subversion will not allow you to relocate a subtree of the working
copy. If you\textquotesingle re going to relocate the working copy at
all, you must relocate the whole thing. This is to protect the integrity
of the working copy metadata and behavior as a whole. (And really,
you\textquotesingle d be hard pressed to come up with a compelling
reason to relocate only a piece of your working copy anyway.)

Let\textquotesingle s look at one final relocation opportunity. After
using HTTP access for some time, the company moves to SSL-only access.
Now, the only change to the repository URL is that the scheme goes from
being \texttt{http://} to being \texttt{https://}. There are two
different ways that we could make our working copy reflect ths change.
The first is to do exactly as we\textquotesingle ve done before and
relocate to the new repository URL.

\begin{verbatim}
$ svn relocate https://svn.company.com/repos/trunk
$
\end{verbatim}

But we have another option here, too. We could simply ask Subversion to
swap out the changed prefixes of the URL.

\begin{verbatim}
$ svn relocate http https
$
\end{verbatim}

Either approach leaves us a working copy whose metadata has been updated
to point to the right repository location.

By default, \texttt{svn\ relocate} will traverse any external working
copies nested within your working copy and attempt relocation of those
working copies, too. Use the \texttt{-\/-ignore-externals} option to
disable this behavior.

\section{svn resolve}\label{svn.ref.svn.c.resolve}

\emph{Resolve conflicts on working copy files or directories.}

\textbf{Synopsis}

\texttt{svn\ resolve\ {[}PATH...{]}}

\subsection{Description}

Resolve ``conflicted'' state on working copy files or directories. This
routine does not semantically resolve conflict markers; however, it
replaces the conflicted item with the version specified (interactively
or via the \texttt{-\/-accept} argument) and then removes
conflict-related artifact files. This allows
\textless PATH\textgreater{} to be committed again---that is, it tells
Subversion that the conflicts have been ``resolved.''

See \hyperref[svn.tour.cycle.resolve]{Resolve Any Conflicts} for an
in-depth look at resolving conflicts.

\subsection{Options}

\begin{verbatim}
--accept ACTION
--depth ARG
--quiet (-q)
--recursive (-R)
--targets FILENAME
\end{verbatim}

\subsection{Examples}

Here\textquotesingle s an example where, after a postponed conflict
resolution during update, \texttt{svn\ resolve} replaces the all
conflicts in file \texttt{foo.c} with your edits:

\begin{verbatim}
$ svn update
Updating '.':
Conflict discovered in 'foo.c'.
Select: (p) postpone, (df) diff-full, (e) edit,
        (mc) mine-conflict, (tc) theirs-conflict,
        (s) show all options: p
C    foo.c
Updated to revision 5.
Summary of conflicts:
  Text conflicts: 1
$ svn resolve --accept mine-full foo.c
Resolved conflicted state of 'foo.c'
$
\end{verbatim}

\section{svn resolved}\label{svn.ref.svn.c.resolved}

\emph{\emph{Deprecated}. Remove ``conflicted'' state on working copy
files or directories.}

\textbf{Synopsis}

\texttt{svn\ resolved\ PATH...}

\subsection{Description}

This command has been deprecated in favor of running
\texttt{svn\ resolve\ -\/-accept\ working\ PATH}. See
\hyperref[svn.ref.svn.c.resolve]{svn resolve} in the preceding section
for details.

Remove ``conflicted'' state on working copy files or directories. This
routine does not semantically resolve conflict markers; it merely
removes conflict-related artifact files and allows
\textless PATH\textgreater{} to be committed again; that is, it tells
Subversion that the conflicts have been ``resolved.'' See
\hyperref[svn.tour.cycle.resolve]{Resolve Any Conflicts} for an in-depth
look at resolving conflicts.

\subsection{Options}

\begin{verbatim}
--depth ARG
--quiet (-q)
--recursive (-R)
--targets FILENAME
\end{verbatim}

\subsection{Examples}

If you get a conflict on an update, your working copy will sprout three
new files:

\begin{verbatim}
$ svn update
Updating '.':
C    foo.c
Updated to revision 31.
Summary of conflicts:
  Text conflicts: 1
$ ls foo.c*
foo.c
foo.c.mine
foo.c.r30
foo.c.r31
$
\end{verbatim}

Once you\textquotesingle ve resolved the conflict and \texttt{foo.c} is
ready to be committed, run \texttt{svn\ resolved} to let your working
copy know you\textquotesingle ve taken care of everything.

You \emph{can} just remove the conflict files and commit, but
\texttt{svn\ resolved} fixes up some bookkeeping data in the working
copy administrative area in addition to removing the conflict files, so
we recommend that you use this command.

\section{svn revert}\label{svn.ref.svn.c.revert}

\emph{Undo all local edits.}

\textbf{Synopsis}

\texttt{svn\ revert\ PATH...}

\subsection{Description}

Reverts any local changes to a file or directory and resolves any
conflicted states. \texttt{svn\ revert} will revert not only the
contents of an item in your working copy, but also any property changes.
Finally, you can use it to undo any scheduling operations that you may
have performed (e.g., files scheduled for addition or deletion can be
``unscheduled'').

\subsection{Options}

\begin{verbatim}
--changelist (--cl) ARG
--depth ARG
--quiet (-q)
--recursive (-R)
--targets FILENAME
\end{verbatim}

\subsection{Examples}

Discard changes to a file:

\begin{verbatim}
$ svn revert foo.c
Reverted foo.c
\end{verbatim}

If you want to revert a whole directory of files, use the
\texttt{-\/-depth=infinity} option:

\begin{verbatim}
$ svn revert --depth=infinity .
Reverted newdir/afile
Reverted foo.c
Reverted bar.txt
\end{verbatim}

Lastly, you can undo any scheduling operations:

\begin{verbatim}
$ svn add mistake.txt whoops
A         mistake.txt
A         whoops
A         whoops/oopsie.c

$ svn revert mistake.txt whoops
Reverted mistake.txt
Reverted whoops

$ svn status
?       mistake.txt
?       whoops
\end{verbatim}

\texttt{svn\ revert} is inherently dangerous, since its entire purpose
is to throw away data---namely, your uncommitted changes. Once
you\textquotesingle ve reverted, Subversion provides \emph{no way} to
get back those uncommitted changes.

If you provide no targets to \texttt{svn\ revert}, it will do nothing.
To protect you from accidentally losing changes in your working copy,
\texttt{svn\ revert} requires you to explicitly provide at least one
target.

\section{svn status (stat, st)}\label{svn.ref.svn.c.status}

\emph{Print the status of working copy files and directories.}

\textbf{Synopsis}

\texttt{svn\ status\ {[}PATH...{]}}

\subsection{Description}

Print the status of working copy files and directories. With no
arguments, it prints only locally modified items (no repository access).
With \texttt{-\/-show-updates} (\texttt{-u}), it adds working revision
and server out-of-date information. With \texttt{-\/-verbose}
(\texttt{-v}), it prints full revision information on every item. With
\texttt{-\/-quiet} (\texttt{-q}), it prints only summary information
about locally modified items.

The first seven columns in the output are each one character wide, and
each column gives you information about a different aspect of each
working copy item.

The first column indicates that an item was added, deleted, or otherwise
changed:

\begin{description}
\item[\texttt{\textquotesingle{}\ \textquotesingle{}}]
No modifications.
\item[\texttt{\textquotesingle{}A\textquotesingle{}}]
Item is scheduled for addition.
\item[\texttt{\textquotesingle{}D\textquotesingle{}}]
Item is scheduled for deletion.
\item[\texttt{\textquotesingle{}M\textquotesingle{}}]
Item has been modified.
\item[\texttt{\textquotesingle{}R\textquotesingle{}}]
Item has been replaced in your working copy. This means the file was
scheduled for deletion, and then a new file with the same name was
scheduled for addition in its place.
\item[\texttt{\textquotesingle{}C\textquotesingle{}}]
The contents (as opposed to the properties) of the item conflict with
updates received from the repository.
\item[\texttt{\textquotesingle{}X\textquotesingle{}}]
Item is present because of an externals definition.
\item[\texttt{\textquotesingle{}I\textquotesingle{}}]
Item is being ignored (e.g., with the \texttt{svn:ignore} property).
\item[\texttt{\textquotesingle{}?\textquotesingle{}}]
Item is not under version control.
\item[\texttt{\textquotesingle{}!\textquotesingle{}}]
Item is missing (e.g., you moved or deleted it without using
\texttt{svn}). This also indicates that a directory is incomplete (a
checkout or update was interrupted).
\item[\texttt{\textquotesingle{}\textasciitilde{}\textquotesingle{}}]
Item is versioned as one kind of object (file, directory, link), but has
been replaced by a different kind of object.
\end{description}

The second column tells the status of a file\textquotesingle s or
directory\textquotesingle s properties:

\begin{description}
\item[\texttt{\textquotesingle{}\ \textquotesingle{}}]
No modifications.
\item[\texttt{\textquotesingle{}M\textquotesingle{}}]
Properties for this item have been modified.
\item[\texttt{\textquotesingle{}C\textquotesingle{}}]
Properties for this item are in conflict with property updates received
from the repository.
\end{description}

The third column is populated only if the working copy directory is
locked (see \hyperref[svn.tour.cleanup]{Sometimes You Just Need to Clean
Up}):

\begin{description}
\item[\texttt{\textquotesingle{}\ \textquotesingle{}}]
Item is not locked.
\item[\texttt{\textquotesingle{}L\textquotesingle{}}]
Item is locked.
\end{description}

The fourth column is populated only if the item is scheduled for
addition-with-history:

\begin{description}
\item[\texttt{\textquotesingle{}\ \textquotesingle{}}]
No history scheduled with commit.
\item[\texttt{\textquotesingle{}+\textquotesingle{}}]
History scheduled with commit.
\end{description}

The fifth column is populated only if the item is switched relative to
its parent (see \hyperref[svn.branchmerge.switchwc]{Traversing
Branches}):

\begin{description}
\item[\texttt{\textquotesingle{}\ \textquotesingle{}}]
Item is a child of its parent directory.
\item[\texttt{\textquotesingle{}S\textquotesingle{}}]
Item is switched.
\end{description}

The sixth column is populated with lock information:

\begin{description}
\item[\texttt{\textquotesingle{}\ \textquotesingle{}}]
When \texttt{-\/-show-updates} (\texttt{-u}) is used, this means the
file is not locked. If \texttt{-\/-show-updates} (\texttt{-u}) is
\emph{not} used, this merely means that the file is not locked in this
working copy.
\item[\texttt{\textquotesingle{}K\textquotesingle{}}]
File is locked in this working copy.
\item[\texttt{\textquotesingle{}O\textquotesingle{}}]
File is locked either by another user or in another working copy. This
appears only when \texttt{-\/-show-updates} (\texttt{-u}) is used.
\item[\texttt{\textquotesingle{}T\textquotesingle{}}]
File was locked in this working copy, but the lock has been ``stolen''
and is invalid. The file is currently locked in the repository. This
appears only when \texttt{-\/-show-updates} (\texttt{-u}) is used.
\item[\texttt{\textquotesingle{}B\textquotesingle{}}]
File was locked in this working copy, but the lock has been ``broken''
and is invalid. The file is no longer locked. This appears only when
\texttt{-\/-show-updates} (\texttt{-u}) is used.
\end{description}

The seventh column is populated only if the item is the victim of a tree
conflict:

\begin{description}
\item[\texttt{\textquotesingle{}\ \textquotesingle{}}]
Item is not the victim of a tree conflict.
\item[\texttt{\textquotesingle{}C\textquotesingle{}}]
Item is the victim of a tree conflict.
\end{description}

The eighth column is always blank.

The out-of-date information appears in the ninth column (only if you
pass the \texttt{-\/-show-updates} (\texttt{-u}) option):

\begin{description}
\item[\texttt{\textquotesingle{}\ \textquotesingle{}}]
The item in your working copy is up to date.
\item[\texttt{\textquotesingle{}*\textquotesingle{}}]
A newer revision of the item exists on the server.
\end{description}

The remaining fields are variable width and delimited by spaces. The
working revision is the next field if the \texttt{-\/-show-updates}
(\texttt{-u}) or \texttt{-\/-verbose} (\texttt{-v}) option is passed.

If the \texttt{-\/-verbose} (\texttt{-v}) option is passed, the last
committed revision and last committed author are displayed next.

The working copy path is always the final field, so it can include
spaces.

\subsection{Options}

\begin{verbatim}
--changelist (--cl) ARG
--depth ARG
--ignore-externals
--incremental
--no-ignore
--quiet (-q)
--show-updates (-u)
--verbose (-v)
--xml
\end{verbatim}

\subsection{Examples}

This is the easiest way to find out what changes you have made to your
working copy:

\begin{verbatim}
$ svn status wc
 M      wc/bar.c
A  +    wc/qax.c
\end{verbatim}

If you want to find out what files in your working copy are out of date,
pass the \texttt{-\/-show-updates} (\texttt{-u}) option (this will
\emph{not} make any changes to your working copy). Here you can see that
\texttt{wc/foo.c} has changed in the repository since we last updated
our working copy:

\begin{verbatim}
$ svn status -u wc
 M            965    wc/bar.c
        *     965    wc/foo.c
A  +          965    wc/qax.c
Status against revision:    981
\end{verbatim}

\texttt{-\/-show-updates} (\texttt{-u}) \emph{only} places an asterisk
next to items that are out of date (i.e., items that will be updated
from the repository if you later use \texttt{svn\ update}).
\texttt{-\/-show-updates} (\texttt{-u}) does \emph{not} cause the status
listing to reflect the repository\textquotesingle s version of the item
(although you can see the revision number in the repository by passing
the \texttt{-\/-verbose} (\texttt{-v}) option).

The most information you can get out of the status subcommand is as
follows:

\begin{verbatim}
$ svn status -u -v wc
 M            965       938 sally        wc/bar.c
        *     965       922 harry        wc/foo.c
A  +          965       687 harry        wc/qax.c
              965       687 harry        wc/zig.c
Status against revision:   981
\end{verbatim}

Lastly, you can get \texttt{svn\ status} output in XML format with the
\texttt{-\/-xml} option:

\begin{verbatim}
$ svn status --xml wc
<?xml version="1.0"?>
<status>
<target
   path="wc">
<entry
   path="qax.c">
<wc-status
   props="none"
   item="added"
   revision="0">
</wc-status>
</entry>
<entry
   path="bar.c">
<wc-status
   props="normal"
   item="modified"
   revision="965">
<commit
   revision="965">
<author>sally</author>
<date>2008-05-28T06:35:53.048870Z</date>
</commit>
</wc-status>
</entry>
</target>
</status>
\end{verbatim}

For many more examples of \texttt{svn\ status}, see
\hyperref[svn.tour.cycle.examine.status]{See an overview of your
changes}.

\section{svn switch (sw)}\label{svn.ref.svn.c.switch}

\emph{Update working copy to a different URL.}

\textbf{Synopsis}

\texttt{svn\ switch\ URL{[}@PEGREV{]}\ {[}PATH{]}}

\texttt{svn\ switch\ -\/-relocate\ FROM\ TO\ {[}PATH...{]}}

\subsection{Description}

The first variant of this subcommand (without the \texttt{-\/-relocate}
option) updates your working copy to point to a new URL. This is the
Subversion way to make a working copy begin tracking a new branch. If
specified, \textless PEGREV\textgreater{} determines in which revision
the target is first looked up. See
\hyperref[svn.branchmerge.switchwc]{Traversing Branches} for an in-depth
look at switching.

Beginning with Subversion 1.7, the \texttt{svn\ switch} command will
demand by default that the URL to which you are switching your working
copy shares a common ancestry with item that the working copy currently
reflects. You can override this behavior by specifying the
\texttt{-\/-ignore-ancestry} option.

If \texttt{-\/-force} is used, unversioned obstructing paths in the
working copy do not automatically cause a failure if the switch attempts
to add the same path. If the obstructing path is the same type (file or
directory) as the corresponding path in the repository, it becomes
versioned but its contents are left untouched in the working copy. This
means that an obstructing directory\textquotesingle s unversioned
children may also obstruct and become versioned. For files, any content
differences between the obstruction and the repository are treated like
a local modification to the working copy. All properties from the
repository are applied to the obstructing path.

As with most subcommands, you can limit the scope of the switch
operation to a particular tree depth using the \texttt{-\/-depth}
option. Alternatively, you can use the \texttt{-\/-set-depth} option to
set a new ``sticky'' working copy depth on the switch target.

The \texttt{-\/-relocate} option is deprecated as of Subversion 1.7. Use
\texttt{svn\ relocate} (described in
\hyperref[svn.ref.svn.c.relocate]{svn relocate}) to perform working copy
relocation instead.

\subsection{Options}

\begin{verbatim}
--accept ACTION
--depth ARG
--diff3-cmd CMD
--force
--ignore-ancestry
--ignore-externals
--quiet (-q)
--relocate
--revision (-r) REV
--set-depth ARG
\end{verbatim}

\subsection{Examples}

If you\textquotesingle re currently inside the directory
\texttt{vendors}, which was branched to \texttt{vendors-with-fix}, and
you\textquotesingle d like to switch your working copy to that branch:

\begin{verbatim}
$ svn switch http://svn.red-bean.com/repos/branches/vendors-with-fix .
U    myproj/foo.txt
U    myproj/bar.txt
U    myproj/baz.c
U    myproj/qux.c
Updated to revision 31.
\end{verbatim}

To switch back, just provide the URL to the location in the repository
from which you originally checked out your working copy:

\begin{verbatim}
$ svn switch http://svn.red-bean.com/repos/trunk/vendors .
U    myproj/foo.txt
U    myproj/bar.txt
U    myproj/baz.c
U    myproj/qux.c
Updated to revision 31.
\end{verbatim}

You \emph{can} switch just part of your working copy to a branch if you
don\textquotesingle t want to switch your entire working copy, but this
is not generally recommended. It\textquotesingle s too easy to forget
that you\textquotesingle ve done this and wind up accidentally making
and committing changes both to the switched and unswitched portions of
your tree.

\section{svn unlock}\label{svn.ref.svn.c.unlock}

\emph{Unlock working copy paths or URLs.}

\textbf{Synopsis}

\texttt{svn\ unlock\ TARGET...}

\subsection{Description}

Unlock each \textless TARGET\textgreater. If any
\textless TARGET\textgreater{} is locked by another user or no valid
lock token exists in the working copy, print a warning and continue
unlocking the rest of the \textless TARGET\textgreater s. Use
\texttt{-\/-force} to break a lock belonging to another user or working
copy.

\subsection{Options}

\begin{verbatim}
--force
--targets FILENAME
\end{verbatim}

\subsection{Examples}

Unlock two files in your working copy:

\begin{verbatim}
$ svn unlock tree.jpg house.jpg
'tree.jpg' unlocked.
'house.jpg' unlocked.
\end{verbatim}

Unlock a file in your working copy that is currently locked by another
user:

\begin{verbatim}
$ svn unlock tree.jpg
svn: E195013: 'tree.jpg' is not locked in this working copy
$ svn unlock --force tree.jpg
'tree.jpg' unlocked.
\end{verbatim}

Unlock a file without a working copy:

\begin{verbatim}
$ svn unlock http://svn.red-bean.com/repos/test/tree.jpg
'tree.jpg unlocked.
\end{verbatim}

For further details, see \hyperref[svn.advanced.locking]{Locking}.

\section{svn update (up)}\label{svn.ref.svn.c.update}

\emph{Update your working copy.}

\textbf{Synopsis}

\texttt{svn\ update\ {[}PATH...{]}}

\subsection{Description}

\texttt{svn\ update} brings changes from the repository into your
working copy. If no revision is given, it brings your working copy up to
date with the \texttt{HEAD} revision. Otherwise, it synchronizes the
working copy to the revision given by the \texttt{-\/-revision}
(\texttt{-r}) option. As part of the synchronization,
\texttt{svn\ update} also removes any stale locks (see
\hyperref[svn.tour.cleanup]{Sometimes You Just Need to Clean Up}) found
in the working copy.

For each updated item, it prints a line that starts with a character
reporting the action taken. These characters have the following meaning:

\begin{description}
\item[\texttt{A}]
Added
\item[\texttt{B}]
Broken lock (third column only)
\item[\texttt{D}]
Deleted
\item[\texttt{U}]
Updated
\item[\texttt{C}]
Conflicted
\item[\texttt{G}]
Merged
\item[\texttt{E}]
Existed
\end{description}

A character in the first column signifies an update to the actual file,
whereas updates to the file\textquotesingle s properties are shown in
the second column. Lock information is printed in the third column.

As with most subcommands, you can limit the scope of the update
operation to a particular tree depth using the \texttt{-\/-depth}
option. Alternatively, you can use the \texttt{-\/-set-depth} option to
set a new ``sticky'' working copy depth on the update target.

\subsection{Options}

\begin{verbatim}
--accept ACTION
--changelist (--cl) ARG
--depth ARG
--diff3-cmd CMD
--editor-cmd CMD
--force
--ignore-externals
--parents
--quiet (-q)
--revision (-r) REV
--set-depth ARG
\end{verbatim}

\subsection{Examples}

Pick up repository changes that have happened since your last update:

\begin{verbatim}
$ svn update
Updating '.':
A    newdir/toggle.c
A    newdir/disclose.c
A    newdir/launch.c
D    newdir/README
Updated to revision 32.
\end{verbatim}

You can also ``update'' your working copy to an older revision
(Subversion doesn\textquotesingle t have the concept of ``sticky'' files
like CVS does):

\begin{verbatim}
$ svn update -r30
Updating '.':
A    newdir/README
D    newdir/toggle.c
D    newdir/disclose.c
D    newdir/launch.c
U    foo.c
Updated to revision 30.
\end{verbatim}

If you want to examine an older revision of a single file, you may want
to use \texttt{svn\ cat} instead---it won\textquotesingle t change your
working copy.

\texttt{svn\ update} is also the primary mechanism used to configure
sparse working copies. When used with the \texttt{-\/-set-depth}, the
update operation will omit or reenlist individual working copy members
by modifying their recorded ambient depth to the depth you specify
(fetching information from the repository as necessary). See
\hyperref[svn.advanced.sparsedirs]{Sparse Directories} for more about
sparse directories.

You can update multiple targets with a single invocation, and Subversion
will not only gracefully skip any unversioned targets you provide it,
but as of Subversion 1.7 will also include a post-update summary of all
the updates it performed:

\begin{verbatim}
$ cd my-projects
$ svn update *
Updating 'calc':
U    button.c
U    integer.c
Updated to revision 394.
Skipped 'tempfile.tmp'
Updating 'paint':
A    palettes.c
U    brushes.c
Updated to revision 60.
Updating 'ziptastic':
At revision 43.
Summary of updates:
  Updated 'calc' to r394.
  Updated 'paint' to r60.
  Updated 'ziptastic' to r43.
Summary of conflicts:
  Skipped paths: 1
$
\end{verbatim}

\section{svn upgrade}\label{svn.ref.svn.c.upgrade}

\emph{Upgrade the metadata storage format for a working copy.}

\textbf{Synopsis}

\texttt{svn\ upgrade\ {[}PATH...{]}}

\subsection{Description}

As new versions of Subversion are released, the format used for the
working copy metadata changes to accomodate new features or fix bugs.
Older versions of Subversion would automatically upgrade working copies
to the new format the first time the working copy was used by the new
version of the software. Beginning with Subversion 1.7, working copy
upgrades must be explicitly performed at the user\textquotesingle s
request. \texttt{svn\ upgrade} is the subcommand used to trigger that
upgrade process.

\subsection{Options}

\begin{verbatim}
--quiet (-q)
\end{verbatim}

\subsection{Examples}

If you attempt to use Subversion 1.7 on a working copy created with an
older version of Subversion, you will see an error:

\begin{verbatim}
$ svn status
svn: E155036: Please see the 'svn upgrade' command
svn: E155036: Working copy '/home/sally/project' is too old (format 10, create
d by Subversion 1.6)
$
\end{verbatim}

Use the \texttt{svn\ upgrade} command to upgrade the working copy to the
most recent metadata format supported by your version of Subversion.

\begin{verbatim}
$ svn upgrade
Upgraded '.'
Upgraded 'A'
Upgraded 'A/B'
Upgraded 'A/B/E'
Upgraded 'A/B/F'
Upgraded 'A/C'
Upgraded 'A/D'
Upgraded 'A/D/G'
Upgraded 'A/D/H'
$ svn status
D       A/B/E/alpha
M       A/D/gamma
A       A/newfile
$
\end{verbatim}

Notice that \texttt{svn\ upgrade} preserved the local modifications
present in the working copy at the time of the upgrade, which were
introduced by the version of Subversion previously used to manipulate
this working copy.

As was the case with automatically upgraded working copies in the past,
explicitly upgraded working copies will be unusable by older versions of
Subversion, too.

